% !TEX program = pdflatex
% !TEX encoding = UTF-8 Unicode

% Plantilla, basada en la clase `scrbook` del paquete KOMA-script,  para la elaboración de un TFG siguiendo las directrices del la comisión del Grado en Matemáticas de la Universidad de Granada.

% Francisco Torralbo Torralbo

\documentclass[print, color]{ugrTFG}

% VERSIÓN ELECTRÓNICA PARA TABLETA
% Cambiando la opción "print" por "tablet" generaremos un pdf adaptado para leerlo en tabletas de 9 pulgadas.

% -------------------------------------------------------------------
% INFORMACIÓN DEL TFG Y EL AUTOR
% -------------------------------------------------------------------

\newcommand{\miTitulo}{Título del trabajo\xspace}
\newcommand{\miNombre}{Isabel María Moreno Cuadrado\xspace}
\newcommand{\miGrado}{Doble grado en Matemáticas e Ingeniería Informática}
\newcommand{\miFacultad}{Facultad de Ciencias y Facultad de la ETSIIT \textbf{(falta logo)}}
\newcommand{\miUniversidad}{Universidad de Granada}

% Añadir tantos tutores como sea necesario separando cada uno de ellos mediante el comando `\medskip` y una línea en blanco
\newcommand{\miTutor}{
  Nombre del tutor 1 \\ \emph{Departamento del tutor 1} 

  % Añadir tantos tutores como sea necesario. 

  \medskip
  Nombre del tutor 2 \\ \emph{Departamento del tutor 2}
}
\newcommand{\miCurso}{2023-2024\xspace}

\hypersetup{
	pdftitle={\miTitulo},
	pdfauthor={\textcopyright\ \miNombre, \miFacultad, \miUniversidad}
}




\usepackage{pgfplots}
\usepackage{amsmath}
\usepackage{amsfonts}
\usepackage{amssymb}
\usepackage{mathtools}
\usepackage{algorithm}
\usepackage{algpseudocode}
\usepackage{mathrsfs}
\usepackage{enumitem}

\begin{document}

\maketitle

% -------------------------------------------------------------------
% FRONTMATTER
% -------------------------------------------------------------------
\frontmatter % Desactiva la numeración de capítulos y usa numeración romana para las páginas

% !TeX root = ../tfg.tex
% !TeX encoding = utf8
%
%*******************************************************
% Declaración de originalidad
%*******************************************************

\thispagestyle{empty}

\hfill\vfill

\textsc{Declaración de originalidad}\\\bigskip

D./Dña. \miNombre \\\medskip

Declaro explícitamente que el trabajo presentado como Trabajo de Fin de Grado (TFG), correspondiente al curso académico \miCurso, es original, entendido esto en el sentido de que no he utilizado para la elaboración del trabajo fuentes sin citarlas debidamente.
\medskip

En Granada a \today 
\vspace{3cm}
\begin{center} 
Fdo: \miNombre 

\end{center}

\vfill

\cleardoublepage
\endinput
   
% !TeX root = ../tfg.tex
% !TeX encoding = utf8

%*******************************************************
% Dedication
%*******************************************************
\thispagestyle{empty}
\phantomsection 
\pdfbookmark[1]{Dedicatoria}{Dedicatoria}

\hfill
\vfill

\begin{flushright}
\itshape
Dedicatoria (opcional) \\
Ver archivo \texttt{preliminares/dedicatoria.tex}
\end{flushright}

\vfill

\cleardoublepage
\endinput
                % Opcional
% !TeX root = ../tfg.tex
% !TeX encoding = utf8

%*******************************************************
% Table of Contents
%*******************************************************
\phantomsection
\pdfbookmark[0]{\contentsname}{toc}

\setcounter{tocdepth}{2} % <-- 2 includes up to subsections in the ToC
\setcounter{secnumdepth}{3} % <-- 3 numbers up to subsubsections

\tableofcontents 

%*******************************************************
% List of Figures and of the Tables
%*******************************************************

    % *******************************************************
    %  List of Figures
    % *******************************************************    
    \phantomsection 
    % \listoffigures

    %*******************************************************
    % List of Tables
    %*******************************************************
    \phantomsection 
    % \listoftables
    
    %*******************************************************
    % List of Listings
    % The package \usepackage{listings} is needed
    %*******************************************************      
	  % \phantomsection 
    % \renewcommand{\lstlistlistingname}{Listados de código}
    % \lstlistoflistings 

\cleardoublepage
            
% !TeX root = ../tfg.tex
% !TeX encoding = utf8

%*******************************************************
% Agradecimientos
%*******************************************************

\chapter{Agradecimientos}

Agradecimientos (opcional, ver archivo \texttt{preliminares/agradecimiento.tex}).

\cleardoublepage
\endinput
            % Opcional

% !TeX root = ../tfg.tex
% !TeX encoding = utf8
%
%*******************************************************
% Summary
%*******************************************************

\selectlanguage{english}
\chapter{Summary}

An english summary of the project (around 800 and 1500 words are recommended).

File: \texttt{preliminares/summary.tex}


% Al finalizar el resumen en inglés, volvemos a seleccionar el idioma español para el documento
\selectlanguage{spanish} 
\endinput
                    
% !TeX root = ../tfg.tex
% !TeX encoding = utf8
%
%*******************************************************
% Introducción
%*******************************************************

% \manualmark
% \markboth{\textsc{Introducción}}{\textsc{Introducción}} 

\chapter{Introducción}

\noindent Al aplicar la transformada de Fourier a una señal en el dominio del tiempo, obtenemos su representación en el dominio de la frecuencia. Esto nos permite examinar la señal original como una composición en términos de componentes sinusoidales.

\noindent Una imagen no es más que una señal visual, y es así susceptible de ser representada en el dominio de la frecuencia. Allí, trabajando con la magnitud y la fase, no solo podremos recuperar la imagen original, sino que ciertas tareas como extracción de características, compresión o  eliminación de ruido, se verán significativamente simplificadas.

\noindent Dado que partimos de datos discretos, y que los ordenadores solo pueden manejar sumas finitas, surge la  imperiosa necesidad de definir en este ambiente una 'aproximación' de la transformada de Fourier.


\noindent \textbf{Completar al final del trabajo.}


\endinput
               

% -------------------------------------------------------------------
% MAINMATTER
% -------------------------------------------------------------------
\mainmatter % activa la numeración de capítulos, resetea la numeración de las páginas y usa números arábigos

\part{Parte Matemática} % Dividir un TFG en partes 
% !TeX root = ../tfg.tex
% !TeX encoding = utf8
%
%*******************************************************
% Introducción
%*******************************************************

% \manualmark
% \markboth{\textsc{Introducción}}{\textsc{Introducción}} 
\chapter{Introducción}
 hablar sobre quien es $L^1(\mathbb{R}^n)$ y definir traslacion y mu  y todo eso  quizas es mejor en otro capitulo ns 
% !TeX root = ../tfg.tex
% !TeX encoding = utf8

\chapter{Transformada de Fourier en $\mathscr{L}^1(\mathbb{R}^n)$}
\section{Definición}

Como se examinará en esta sección, existen múltiples enfoques para establecer la definición de la Transformada de Fourier en $\mathbb{R}^n$. La formulación que adoptaremos se basa en las presentadas en las referencias~\cite{FourierAnalysisApplications} y~\cite{FirstCourseFourier}. No obstante, se considerarán y discutirán otras definiciones alternativas, como las propuestas en~\cite{FourierTransformsClassic} y~\cite{Riviere2023Fourier}.

\begin{definicion}\label{def:definicion}
Sea $f \in \mathscr{L}^1(\mathbb{R}^n)$. La Transformada de Fourier de $f$ es  la función $\widehat{f} : \mathbb{R}^n \rightarrow \mathbb{C}$ definida por 
\begin{equation}
    \widehat{f}(y) = \int_{\mathbb{R}^n} f(x) e^{-2\pi i \langle x, y \rangle} \, dx \quad \forall y \in \mathbb{R}^n.
\end{equation}
\end{definicion}


\begin{observacion}
    Lo primero que debemos mencionar es que la definición anterior es legítima: dado $y \in \mathbb{R}^n$, la función $x \rightarrow f(x)e^{-2\pi i\langle x, y \rangle}$ es medible por ser el producto de una función medible $f$ y la función continua $x \rightarrow e^{-2 \pi i \langle x, y \rangle}$, y además 
    \begin{equation}
        \int_{\mathbb{R}^n}|f(x)e^{-2\pi i\langle x, y \rangle}| \, dx  = \int_{\mathbb{R}^n}|f(x)| \, dx < \infty,
    \end{equation}
\end{observacion}
\noindent luego es integrable en $\mathbb{R}^n$ para cada $y\in\mathbb{R}^n$.

\begin{observacion} \label{obs2}
Como el lector puede imaginar, uno de nuestros objetivos es ser capaces de reconstruir la función $f$ a partir de $\widehat{f}$. Esta tarea se pretende llevar a cabo mediante la fórmula 
\begin{equation}
    f(x) = \int_{\mathbb{R}^n} \widehat{f}(y) e^{2\pi i \langle x, y \rangle} \, dy .
\end{equation}
Esto implica que la función Transformada \( \widehat{f} \) debe ser integrable sobre \( \mathbb{R}^n \), lo cual puede no ser siempre el caso, incluso para funciones aparentemente simples. En situaciones donde la función Transformada no cumpla con el requisito de ser integrable, se consideran otras interpretaciones alternativas a la fórmula propuesta.
\end{observacion}

\begin{observacion}
Es interesante observar que el concepto de la Transformada de Fourier tiene sentido para funciones definidas casi por doquier en $\mathbb{R}^n$.
Además, si $f=g$ c.p.d. en $\mathbb{R}^n$, entonces se tiene que $\widehat{f}=\widehat{g}$. Por tanto, la Transformada de Fourier se transfiere naturalmente al espacio $L^1(\mathbb{R}^n)$.

\end{observacion}

\begin{observacion}
Es frecuente encontrar otros textos en los que la transformada se define mediante expresiones diferentes, como puede ser:
\begin{equation}\label{eq:trans2}
    \int_{\mathbb{R}^n} f(x) e^{-i \langle x, y \rangle} \, dx,
\end{equation}
definición que aparece en~\cite{FourierTransformsClassic}, y es aparentemente más sencilla. Lo cierto es que esto obliga a que la reconstrucción se haga mediante la fórmula 
\begin{equation}
   \frac{1}{(2 \pi)^n} \int_{\mathbb{R}^n} f(x) e^{i \langle x, y \rangle} \, dy ,
\end{equation}
por lo que la constante (ligada  a $2 \pi$) aparece en la reconstrucción, esta vez dependiente del número de dimensiones.

\noindent En general, si definimos la Transformada de Fourier como 
\begin{equation}
    a\int_{\mathbb{R}^n} f(x) e^{-i \langle x, y \rangle} \, dx ,
\end{equation}
esto obliga a que la reconstrucción venga dada por 
\begin{equation}
   b\int_{\mathbb{R}^n} f(x) e^{i \langle x, y \rangle} \, dy, \quad 
\end{equation}
donde $a \cdot b = \frac{1}{(2\pi)^n}$.

\noindent Por tanto, no podemos eludir la presencia de la constante $2 \pi$ en ninguna instancia. Sin embargo, parece que con nuestra elección minimizamos el riesgo de cometer un error, ya que al incorporar la constante en la exponencial, obtenemos fórmulas completamente simétricas y válidas para cualquier dimensión. Pero si el lector aún no ha quedado convencido, es crucial señalar que, usando la definición (\ref{def:definicion}), el Teorema de Convolución (que será el núcleo central de este trabajo) no involucrará la presencia de constantes añadidas (aunque estas constantes sí aparecerán en el contexto de la Transformada  de Fourier y la derivación).
\end{observacion}

\section{Propiedades de la Transformada de Fourier}
Demostraremos algunas propiedades de la Transformada de Fourier para comenzar a familiarizarnos con la definición.

\begin{proposicion}\label{prop_lineal}
    Sean $f,g \in \mathscr{L}^1(\mathbb{R}^n)$ y sean $\alpha,\beta \in \mathbb{C}$. Entonces
    \begin{equation}
         \widehat{\alpha f+\beta g} = \alpha \widehat{f} + \beta \widehat{g}.
    \end{equation}
\end{proposicion}

\begin{proof}
Para cada $y \in \mathbb{R}^n$, se verifica que
\begin{align*}
\widehat{ \alpha f+ \beta g} &= \int_{ \mathbb{R}^n}(\alpha f(x) + \beta g(x)) e^{-2\pi i \langle x, y \rangle} \, dx \\
&= \alpha \int_{\mathbb{R}^n} f(x)e^{-2 \pi i \langle x, y \rangle} \, dx + \beta \int_{\mathbb{R}^n} g(x)e^{-2 \pi i \langle x, y \rangle} \, dx
= \alpha \widehat{f} + \beta \widehat{g}. \qedhere
\end{align*}
\end{proof}

\begin{proposicion}\label{conju}
Sean $f,g \in \mathscr{L}^1(\mathbb{R}^n)$. Entonces 
%\renewcommand{\labelenumi}{\roman{enumi})}

\begin{enumerate}
    \item $\widehat{\overline{f}}(y) = \overline{\widehat{f}(-y)}$,
    \item $\widehat{\widetilde{f}(}y) = \widehat{f}(-y)$,
    \item $\widehat{\overline{\widetilde{f}(}}y) = \overline{\widehat{f}(y)}$.
\end{enumerate}
\end{proposicion}

\begin{proof}
\begin{align*}
      \widehat{\overline{f}}(y) &= \int_{\mathbb{R}^n} \overline{f(x)}e^{-2\pi i \langle x, y \rangle} \, dx = \int_{\mathbb{R}^n} \overline{f(x)e^{2\pi i \langle x, y \rangle}} \, dx =  \overline{\int_{\mathbb{R}^n} f(x)e^{2\pi i \langle x, y \rangle} \, dx } =\overline{\widehat{f}(-y)}, \\[0.2cm]
       \widehat{\widetilde{f}(}y) &= \int_{\mathbb{R}^n} f(-x)e^{-2\pi i \langle x, y \rangle} \, dx = \int_{\mathbb{R}^n} f(x)e^{-2\pi i \langle (-x), y \rangle} \, dx  =  \int_{\mathbb{R}^n} f(x)e^{2\pi i \langle x, (-y) \rangle} \, dx  =\widehat{f}(-y),\\[0.2cm]
       \widehat{\overline{\widetilde{f}(}}y) &= \int_{\mathbb{R}^n} \overline{f(-x)}e^{-2\pi i \langle x, y \rangle} \, dx = \int_{\mathbb{R}^n} \overline{f(-x)e^{2\pi i \langle x, y \rangle}} \, dx  =  \overline{\int_{\mathbb{R}^n} f(x)e^{-2\pi i \langle x, y \rangle} \, dx } =\overline{\widehat{f}(y)}. \qedhere
\end{align*} 
\end{proof}




\noindent Las siguientes proposiciones son útiles en lo que se refiere al cálculo de Transformadas de Fourier. Las demostraciones se reducen a efectuar un cambio de variable adecuado en la integral.

\begin{proposicion}\label{esca}
    Sea $f \in \mathscr{L}^1(\mathbb{R}^n)$ y sea $a \in \mathbb{R}^+$. Entonces, para cada $y \in \mathbb{R}^n$, se tiene:
    \begin{equation}
        \widehat{H_af}(y) = a^{-n}\widehat{f}(a^{-1}y).
    \end{equation}
\end{proposicion}

\begin{proof}
\begin{equation*}
     \widehat{H_af}(y) = \int_{\mathbb{R}^n}f(ax) e^{-2 \pi i \langle x, y \rangle} \, dx = \frac{1}{a^n}\int_{\mathbb{R}^n}f(x) e^{-2 \pi i \langle x, \frac{y}{a} \rangle} \, dx =  a^{-n}\widehat{f}(a^{-1}y). \qedhere
\end{equation*}
\end{proof}


\begin{proposicion}\label{prop:tras}
    Sea $f \in \mathscr{L}^1(\mathbb{R}^n)$ y sea $t \in \mathbb{R}^n$. Entonces, para cada $y \in \mathbb{R}^n$, se tiene:
    \begin{equation}
        \widehat{\tau_tf}(y) =  e^{-2 \pi i \langle t, y \rangle}\widehat{f}(y).
    \end{equation}
    
\end{proposicion}

\begin{proof}
\begin{align*}
        \widehat{\tau_tf}(y) &= \int_{\mathbb{R}^n}f(x-t) e^{-2 \pi i \langle x, y \rangle} \, dx = \int_{\mathbb{R}^n}f(x) e^{-2 \pi i \langle (x+t), y \rangle} \, dx \\ &=  \int_{\mathbb{R}^n}f(x) e^{-2 \pi i \langle x, y \rangle} e^{-2 \pi i \langle t, y \rangle} \, dx = e^{-2 \pi i \langle t, y \rangle}\widehat{f}(y). \qedhere
\end{align*}
\qedhere
\end{proof}


\begin{proposicion}
    Sea $f \in \mathscr{L}^1(\mathbb{R}^n)$ y sea $t \in \mathbb{R}^n$. Entonces, para cada $y \in \mathbb{R}^n$, se tiene:
    \begin{equation*}
        \widehat{\mu_tf}(y) = \widehat{f}(y-t). \qedhere
    \end{equation*}
\end{proposicion}

\begin{proof}
\begin{alignat}{3}
        \widehat{\mu_tf}(y) &= \int_{\mathbb{R}^n}e^{2 \pi i \langle t, x \rangle}f(x) e^{-2 \pi i \langle x, y \rangle} \, dx = \int_{\mathbb{R}^n}f(x) e^{-2 \pi i \langle x, (y-t) \rangle} \, dx   = \widehat{f}(y-t). \qedhere
\end{alignat}
\end{proof}

\begin{proposicion}\label{prop:div}
    Si $f \in \mathscr{L}^1(\mathbb{R}^n)$ 
 es tal que $f(x_1, \ldots, x_n) = f_1(x_1) \cdot f_2(x_2) \cdots f_n(x_n)  
 \\ \quad \forall (x_1,x_2, \ldots, x_n) \in \mathbb{R}^n$, con $f_1,f_2, \ldots, f_n \in \mathscr{L}^1(\mathbb{R}^n)$. Entonces
    \begin{equation}
        \widehat{f}(y_1,y_2 \ldots, y_n) =  \widehat{f_1}(y_1) \cdot \widehat{f_2}(y_2) \cdots \widehat{f_n}(y_n) \quad \forall y = (y_1,y_2, \ldots, y_n) \in \mathbb{R}^n.
    \end{equation}
\end{proposicion}


\begin{proof}
Sea $y = (y_1,y_2, \ldots, y_n) \in \mathbb{R}^n$. Entonces, se tiene que
\begin{equation*}
\begin{split}
    \widehat{f}(y) &= \int_{\mathbb{R}^n} (f_1(x_1)e^{-2\pi i (x_1y_1)} \cdot f_2(x_2)e^{-2\pi i (x_2y_2)} \cdots f_n(x_n)e^{-2\pi i (x_ny_n)} )  \, dx_1dx_2\ldots dx_n \\
    &= \prod_{j=1}^{n}\int_{\mathbb{R}}f_j(x_j)e^{-2\pi i x_j y_j} dx_j =  \prod_{j=1}^{n}\widehat{f_j}(y_j) . \qedhere
\end{split}
\end{equation*}

\end{proof}
\begin{observacion}
    Esta proposición nos permite reducir el cálculo de la Transformada de Fourier de una función integrable en $\mathbb{R}^n$ al caso $n=1$ en determinadas situaciones.
\end{observacion}



\section{Permutación integratoria de la Transformada}
El siguiente teorema nos permite intercambiar la Transformada dentro de la integral. Esta técnica resulta útil como estrategia para abordar ciertas demostraciones.
\begin{teorema}\label{teo:cambio}
    Sean $f,g \in \mathscr{L}^1(\mathbb{R}^n)$. Entonces,
    \begin{equation}
        \int_{\mathbb{R}^n} \widehat{f}(x)g(x) \, dx =  \int_{\mathbb{R}^n} f(x) \widehat{g}(x) \, dx.
    \end{equation}
\end{teorema}

\begin{proof}
Por un lado tenemos que la función $(x,y) \mapsto f(x)g(y)e^{-2 \pi i \langle x, y \rangle}$ es medible, y por otro se tiene que 
\begin{equation}
\begin{aligned}
    &\int_{\mathbb{R}^n}\int_{\mathbb{R}^n}|f(x)g(y)e^{-2 \pi i \langle x, y \rangle}| \,dx \, dy = \\
    &= \int_{\mathbb{R}^n}\int_{\mathbb{R}^n}|f(x)g(y)| \,dx \, dy = \int_{\mathbb{R}^n}|f(x)| \,dx \int_{\mathbb{R}^n}|g(y)| \, dy < \infty,
\end{aligned}
\end{equation}


\noindent Por lo que el Teorema de Fubini-Tonelli asegura que la función es integrable en $\mathbb{R}^{n+n}$ y

\begin{equation}
\begin{aligned}
    \int_{\mathbb{R}^n} \widehat{f}(x)g(x) \, dx &= \int_{\mathbb{R}^n}\int_{\mathbb{R}^n}f(x)g(y)e^{-2 \pi i \langle x, y \rangle} \,dy \,dx  \\
    &= \int_{\mathbb{R}^n}\int_{\mathbb{R}^n}f(x)g(y)e^{-2 \pi i \langle x, y \rangle} \,dx \,dy = \int_{\mathbb{R}^n} f(x) \widehat{g}(x) \, dx. \qedhere
\end{aligned}
\end{equation}
\end{proof}

\section{Magnitud de la Transformada de Fourier}
En esta sección estudiaremos la magnitud de la Transformada de Fourier, proporcionaremos una acotación de esta y demostraremos el Lema de Riemann-Lebesgue.


\begin{teorema}\label{2.2}
    Sea  $f \in \mathscr{L}^1(\mathbb{R}^n)$. Entonces
    \begin{equation}
        || \widehat{f}||_{L^{\infty}} \leq ||f||_{L^{1}}.
    \end{equation}    
\end{teorema}

\begin{proof}
    Para cada $y \in \mathbb{R}^n$, se tiene que
    \begin{equation}
     |\widehat{f}(y)| = \left| \int_{\mathbb{R}^n} f(x) e^{-2\pi i \langle x, y \rangle} \, dx \right| \leq  \int_{\mathbb{R}^n} |f(x) e^{-2\pi i \langle x, y \rangle}| \, dx  =\int_{\mathbb{R}^n} |f(x)| \, dx = ||f||_{1}.
    \end{equation}  
    De aquí se deduce que 
    $      || \widehat{f}||_{L^{\infty}} = \underset{\substack{y \in \mathbb{R}^n}}{\sup}|\widehat{f}(y)| \leq||f||_{1}$.
\end{proof}





\begin{teorema}[Lema de Riemann-Lebesgue]\label{2.3} Sea $f \in \mathscr{L}^1(\mathbb{R}^n)$. Entonces,
\begin{equation}
     \underset{\substack{||y|| \rightarrow \infty}}{\lim}\widehat{f}(y)=0.
\end{equation}
\end{teorema}

\noindent Antes de realizar la prueba recordamos algunas definiciones y resultados conocidos.
\begin{definicion}
Un  intervalo en $\mathbb{R}^n$ será un producto cartesiano de intervalos en $\mathbb{R}$ y denotaremos por $\mathcal{J}$ al  conjunto de todos los intervalos acotados en $\mathbb{R}^n$.
\end{definicion}

\begin{definicion}\label{def:escalonada}
    Llamaremos función escalonada a toda combinación lineal de funciones características de
intervalos acotados, es decir, a toda función $h: \mathbb{R}^n \rightarrow \mathbb{C}$ de la forma
\begin{equation}
h = \sum_{k=1}^{n} \alpha_k \upchi_{J_k} \quad \text{donde} \quad n \in \mathbb{N}, \quad \alpha_1, \alpha_2, \ldots, \alpha_n \in \mathbb{C}, \quad J_1, J_2, \ldots, J_n \in \mathcal{J}.
\end{equation}

\end{definicion}

\begin{observacion}

El conjunto formado por todas las clases de equivalencia que contienen una función escalonada en $\mathbb{R}^n$ es denso en $L^1(\mathbb{R}^n)$. Como consecuencia, para cada $f \in \mathscr{L}^1(\mathbb{R}^n)$, existe una sucesión $\{g_n\}$ de funciones escalonadas que verifica:
\begin{align*}
&\{g_n(x)\} \rightarrow f(x) \text{ p.c.t. } x \in \mathbb{R}^n,
&\lim_{n \to \infty} \int_{\mathbb{R}^n} |f - g_n| \, dx = 0.
\end{align*}
\end{observacion}


\begin{proof}[Demostración del Lema de Riemann-Lebesgue] La prueba está divida en dos partes: primero se demuestra el teorema para cualquier función escalonada, y posteriormente, usando un argumento de densidad, se extiende el resultado a cualquier función integrable en $\mathbb{R}^n$.
\vspace{0.1cm}



\noindent Dado un intervalo $\left[a,b\right]$  en $\mathbb{R}$, consideramos la función característica $\upchi_{\left[a,b\right]}$. Calculamos la Transformada de dicha función.
\begin{equation}
\widehat{\upchi_{\left[a,b\right]}}(0) = \int_{\mathbb{R}} \upchi_{\left[a,b\right]}\, dx = b-a,
\end{equation}
\begin{equation}\label{eq:caract}
\widehat{\upchi_{\left[a,b\right]}}(y) = \int_{\mathbb{R}} \upchi_{\left[a,b\right]}e^{-2\pi i xy} \, dx =\int_{a}^{b} e^{-2\pi i  x y } \, dx = \frac{e^{-2 \pi i y a}-e^{-2 \pi i y b}}{2 \pi i y} \, \, \, \,\forall y \neq 0 .
\end{equation}

\noindent Observamos que
\begin{equation}   
\underset{\substack{|y| \rightarrow \infty}}{\lim}\widehat{\upchi_{\left[a,b\right]}}(y)=0.
\end{equation}


\noindent Tomamos ahora un  intervalo en $\mathbb{R}^n$, $ J = [a_1, b_1] \times [a_2, b_2] \times \cdots \times [a_n, b_n]$ y  la función  característica $\upchi_{J}$. Procedemos calculando la Transformada de Fourier de $\upchi_{J}\in \mathscr{L}^1(\mathbb{R}^n)$:
\begin{align}
    \widehat{\upchi_{J}}(y_1, \ldots, y_n) &= \int_{\mathbb{R}^n} \upchi_{J} e^{-2\pi i \langle x, y \rangle} \, dx \nonumber \\
    &= \int_{a_n}^{b_n}\ \cdots \int_{a_2}^{b_2}\left( \int_{a_1}^{b_1}  e^{-2\pi i x_1y_1 } \, dx_1\right) \, e^{-2\pi i (x_2y_2 + \ldots + x_ny_n) } \, dx_2 \ldots dx_n \\ &= \prod_{j=1}^{n}\left( \int_{a_j}^{b_j}e^{-2 \pi i x_jy_j} dx_j\right) .
\end{align}



\noindent Sea  un elemento $y = (y_1, y_2, \ldots ,y_n)\neq 0 \in \mathbb{R}^n $. Existe $j_0 \in \{1, \ldots, n\}$ tal que $|y_{j_0}| > \frac{||y||}{\sqrt{n}}$. 
Notemos que, para todo $j \in \{1, \ldots, n\}$, se tiene 
\begin{equation}
    \left| \int_{a_j}^{b_j} e^{-2 \pi i x y_j} dx \right| \leq  \int_{a_j}^{b_j} \left| e^{-2 \pi i x y_j} \right| dx  = (b_j-a_j).
\end{equation}

\noindent Por otro lado, 

\begin{equation}
\left|\int_{a_{j_0}}^{b_{j_0}}  e^{-2 \pi i x y_j}  dx  \right| = \left|\frac{e^{-2 \pi i y_{j_0} a_{j_0} }-e^{-2 \pi i y_{j_0}  b_{j_0} }}{2 \pi i _{j_0} }\right| \leq \frac{2}{|2 \pi i y_{j_0}|} <  \frac{\sqrt{n}}{ \pi ||y||}.
\end{equation}

\noindent Por tanto, tenemos que 
\begin{equation}
|\widehat{\upchi_{J}}(y)| \leq \prod_{\substack{j = 1\\ \substack{j \neq j_0}}}^{n}(b_j-a_j)  \cdot \frac{\sqrt{n}}{ \pi ||y||} \leq \max_{1 \leq j_0 \leq n}\prod_{\substack{j = 1\\ \substack{j \neq j_0}}}^{n}(b_j-a_j)  \cdot \frac{\sqrt{n}}{ \pi ||y||} =  c\cdot \frac{\sqrt{n}}{ \pi ||y||},
\end{equation}
donde la constante $c$ no depende de $j_0$.
Atendiendo a esta última expresión, deducimos que 
\begin{equation}   
\underset{\substack{||y|| \rightarrow \infty}}{\lim}\widehat{\upchi_{J}}(y)=0.
\end{equation}

\noindent Dada una función escalonada $g : \mathbb{R}^n \rightarrow \mathbb{R}$, esta será combinación lineal finita de funciones del tipo $X_J$. Podemos concluir 
entonces que, gracias a la aditividad de los límites, el teorema se cumple para $g$. Así, se tendrá que 
\begin{equation}   
\underset{\substack{||y|| \rightarrow \infty}}{\lim}\widehat{g}(y)=0.
\end{equation}
Generalizamos ya el resultado. Dada ahora una función $f \in \mathscr{L}^1(\mathbb{R}^n)$ y $\epsilon > 0$, existe una función escalonada $h \in  \mathscr{L}^1(\mathbb{R}^n)$  de manera que $||f\,-\,h||_{L^1} < \frac{\epsilon}{2}$. Por tanto, por lo probado anteriormente, existe una constante positiva $M$ tal que si $|y| > M$ entonces $|\widehat{h}(y)| <  \frac{\epsilon}{2}$. Luego 
\begin{equation}
    |\widehat{f}(y)| \leq |\widehat{f}(y)-\widehat{h}(y)|+|\widehat{h}(y)| \leq ||f-h||_{L^1} +\frac{\epsilon}{2} \leq \epsilon,
\end{equation}

\noindent para $||y||> M$, de donde se deduce que $\underset{\substack{||y||\rightarrow \infty}}{\lim}\widehat{f}(y)=0$, y hemos terminado.
\end{proof}


\section{Continuidad de la Transformada de Fourier}
A continuación se describe la continuidad de la Transformada de Fourier $\widehat{f}$ dada $f \in  \mathscr{L}^1(\mathbb{R}^n)$.

\begin{teorema}\label{2.4}
Sea $f \in  \mathscr{L}^1(\mathbb{R}^n)$. Entonces la función $\widehat{f}$ es continua en $\mathbb{R}^n$.  
\end{teorema}


\begin{proof}
Podemos tratar la integral que aparece en \ref{def:definicion} como una integral dependiente de un parámetro.
\noindent Sea $\gamma : \mathbb{R}^n \times \mathbb{R}^n \rightarrow \mathbb{C}$ la función dada por
\begin{equation}
    \gamma(x,y) = 
     f(x) e^{-2\pi i \langle x, y \rangle} .
\end{equation}
Se cumple:
\begin{enumerate}
    \item Para cada $y \in \mathbb{R}^n$, la función $x \mapsto \gamma(x,y)$  es integrable en $\mathbb{R}^n$,
    \item Para cada $x \in \mathbb{R}^n$, la función $y \mapsto \gamma(x,y)$  es continua en $\mathbb{R}^n$,
    \item $|\gamma(x,y)| = |f(x) e^{-2\pi i \langle x, y \rangle}| = |f(x)| \quad \forall x,y \in \mathbb{R}^{n}$.
\end{enumerate}
Puesto que $|f|$ es integrable en $\mathbb{R}^n$, el Teorema de continuidad de las integrales dependientes de un parámetro garantiza que la función $\widehat{f}$ es continua en $\mathbb{R}^n$.
\end{proof}


\begin{observacion}
\noindent Los Teoremas \ref{2.4} y \ref{2.3} aseguran que  $\widehat{f}$ es uniformemente continua en $\mathbb{R}^n$. 
\end{observacion}


\begin{observacion}
    Los Teoremas \ref{2.4} y \ref{2.3} aseguran que  $\widehat{f} \in C_0 (\mathbb{R}^n)$, para cada $f \in \mathscr{L}^1( \mathbb{R}^n)$. Teniendo esto en cuenta junto con la Proposición \ref{prop_lineal} y el  Teorema \ref{2.2}, concluimos que la Transformada de Fourier
    \begin{equation}
        \widehat{} :  \mathscr{L}^1(\mathbb{R}^n) \rightarrow C_0(\mathbb{R}^n)
    \end{equation}
    define un operador lineal y continuo.
\end{observacion}







\section{Transformada de Fourier y derivación}
En esta sección se pretende analizar en detalle cómo se interrelacionan la Transformada de Fourier y la operación de derivación. Este análisis conlleva dos enfoques de estudio: uno que supone calcular la Transformada de Fourier de la derivada de una función, y relacionarla con la Transformada de Fourier de la función original; y otro que consiste en abordar la derivación de la función Transformada.


\noindent Comenzaremos por introducir un poco de notación.
\noindent Para $\alpha = (\alpha_1, \alpha_2, \ldots ,\alpha_n) \in (\mathbb{N} \cup {0})^n$ e $y = (y_1,y_2, \ldots, y_n) \in \mathbb{R}^n$ notaremos
\begin{equation}
|\alpha| = \sum_{k=1}^{n}\alpha_k \quad \text{e} \quad y^\alpha = \prod_{k=1}^{n} y_k^{\alpha_k} = y_{1}^{\alpha_1}y_{2}^{\alpha_2}\cdots y_{n}^{\alpha_n}.
\end{equation}

\noindent $\alpha$ se denominará multi-índice.

\noindent Si $f: \mathbb{R}^n \rightarrow \mathbb{C}$ es una función de clase $C^\infty$, escribimos
\begin{equation}
   D_{\alpha}f = \frac{\partial^{|\alpha|f}}{\partial x_{1}^{\alpha_1}\partial x_{2}^{\alpha_2} \cdots \partial x_{n}^{\alpha_n}}.
\end{equation}
\noindent Para $|\alpha| = 0$, decimos que $ D_{\alpha}f = f$.

\begin{definicion}
    Dado $k \in \{1, \ldots, n\}$, denotaremos $e_k$ al vector cuyas componentes son todas nulas salvo la k-ésima. Se tiene que, para $\alpha = e_k$, 
    \begin{equation}
        D_{\alpha}f = \frac{\partial f}{\partial x_k}.
    \end{equation}
\end{definicion}

\begin{definicion}
    Denominaremos $A_{\alpha}$ al conjunto:
    \begin{equation}
        A_{\alpha} = \{\beta \in (\mathbb{N} \cup \{0\})^n :  \, |\beta| \leq |\alpha|\}.
    \end{equation}
\end{definicion}

\subsection{Transformada de la derivada}

\begin{teorema} \label{teoderi}
Dados $f \in  \mathscr{L}^1(\mathbb{R}^n)$ y $k \in \{1,\ldots,n\}$, si existe $ \frac{\partial f}{\partial x_k}$ y además  $ \frac{\partial f}{\partial x_k} \in \mathscr{L}^1(\mathbb{R}^n)$,
entonces
\begin{equation}
    \widehat{\frac{\partial{f}}{\partial{x_k}}}(y) = (2\pi iy_k) \widehat{f}(y)  \quad \forall y \in \mathbb{R}^n.
\end{equation}

\noindent Equivalentemente, para $\alpha = e_k$,
\begin{equation}
    \widehat{D_{\alpha}f}(y) = (2\pi i)^{|\alpha|} y^{\alpha} \widehat{f}(y) \quad \forall y \in \mathbb{R}^n.
\end{equation}
\end{teorema}

\begin{proof}
\noindent

\noindent Por hipótesis sabemos que dado $y \in \mathbb{R}^n$, la función $(x_1,\ldots,x_n) \mapsto  \frac{\partial{f}}{\partial{x_k}}(x_1,\ldots,x_n)e^{- 2 \pi i\sum\limits_{i=1}^{n} x_iy_i}$ es medible  e integrable en $\mathbb{R}^n$:
\begin{equation}
    \left|  \frac{\partial{f}}{\partial{x_k}}(x)e^{- 2 \pi i\sum\limits_{i=1}^{n} x_iy_i}\right| = \left| \frac{\partial{f}}{\partial{x_k}}(x) \right| < \infty.
\end{equation}

\noindent La clave de la demostración radica en usar el Teorema de Fubini y posteriormente aplicar la fórmula de integración por partes en $\mathbb{R}$ a la función $x_k \mapsto \frac{\partial{f}}{\partial{x_k}}(x_1,\ldots,x_n)e^{-2 \pi i x_ky_k}$.
\begin{align}
    \widehat{\frac{\partial{f}}{\partial{x_k}}}(y) 
    &= \int_{\mathbb{R}^n}\frac{\partial{f}}{\partial{x_k}}(x_1,\ldots,x_n)e^{- 2 \pi i\sum\limits_{i=1}^{n} x_iy_i}dx_1\ldots dx_{k-1}dx_kdx_{k+1}\ldots dx_n \\
    &= \int_{\mathbb{R}}\ldots\left(\int_{\mathbb{R}}\frac{\partial{f}}{\partial{x_k}}(x_1,\ldots,x_n)e^{-2 \pi i x_ky_k}dx_k\right) e^{- 2 \pi i\sum\limits_{i=1, i\neq k}^{n} x_iy_i} dx_1 \ldots dx_{k-1}dx_{k+1}\ldots dx_n \\
    &= (2 \pi i y_k)\int_{\mathbb{R}}\ldots\int_{\mathbb{R}} f(x_1,\ldots,x_n)e^{- 2 \pi i\sum\limits_{i=1}^{n} x_iy_i}dx \\
    &= (2\pi iy_k) \widehat{f}(y). \qedhere
\end{align}
\end{proof}




\begin{teorema}
 Sea $\alpha \in (\mathbb{N} \cup \{0\})^n$ y supongamos que existe $D_{\beta}f \, \,  \forall  \beta \in A_{\alpha}$ de modo que
\begin{equation}
    D_{\beta}f \in \mathscr{L}^1(\mathbb{R}^n) \quad \forall  \beta \in A_{\alpha}.
\end{equation}
Entonces
\begin{equation}
    \widehat{D_{\alpha}f}(y) = (2\pi i)^{|\alpha|}y^\alpha \widehat{f}(y) \quad \forall y \in \mathbb{R}^n.
\end{equation}
\end{teorema}

\begin{proof}
    Para un $\alpha \in ( \mathbb{N} \cup \{0\})^n$ cualquiera, basta usar el Teorema \ref{teoderi} y realizar inducción sobre $|\alpha|$. 
\end{proof}




















\subsection{Derivada de la Transformada}
\begin{teorema}\label{teo:der_trans}
   Sea $f \in \mathscr{L}^1(\mathbb{R}^n)$ tal que se cumple la condición 
   \begin{equation}\int_{\mathbb{R}^n}|x_kf(x)| \, dx < \infty,
   \end{equation}
   donde $x_k$ denota la coordenada k-ésima de $x$.
   Entonces $\widehat{f}$ es derivable con respecto a $y_k$, y
   \begin{equation}\label{eq:ecuacion}
    \frac{\partial \widehat{f}}{\partial y_k}(y) = (-2 \pi ix_k f (x))\,\,\widehat{  } \,\, (y) \quad \forall  y \in \mathbb{R}^n.
   \end{equation}
   Equivalentemente, para $\alpha = e_k $,
   \begin{equation}
     D_{\alpha}\widehat{f}(y) = (-2 \pi ix^{\alpha} f (x))\,\,\widehat{  } \,\, (y) \quad \forall  y \in \mathbb{R}^n.
   \end{equation}
\end{teorema}


\begin{proof}
\noindent Sea la función $\gamma : \mathbb{R}^n \times \mathbb{R}^n \rightarrow \mathbb{C}$ definida por 
\begin{equation}
    \gamma(x,y) = 
     f(x) e^{-2\pi i \langle x, y \rangle}.
\end{equation}
Esta función es integrable con respecto la variable $x$ y derivable con respecto a la variable $y_k$.
Tenemos que:
\begin{equation}
    \frac{\partial \gamma}{\partial y_k}(x,y) = \frac{\partial \left(f(x) e^{-2\pi i \langle x, y \rangle}\right)}{\partial y_k} = -2 \pi  i x_k f(x) e^{-2\pi i \langle x, y \rangle}. 
\end{equation}
Además,
\begin{equation}
    \left|\frac{\partial \gamma}{\partial y_k}(x,y)\right| = 2 \pi   |x_k f(x)| \quad \forall x,y \in \mathbb{R}^n.
\end{equation}
Como por hipótesis la función $x_kf \in \mathscr{L}^1(\mathbb{R}^n)$, el Teorema de derivación de integrales dependientes de un parámetro asegura que $\widehat{f}$ es derivable con respecto a $y_k$, y también:
 \begin{align}
    \frac{\partial \widehat{f}}{\partial y_k}(y) &= \int_{\mathbb{R}^n}\frac {\partial}{\partial y_k}\left(f(x) e^{-2\pi i \langle x, y \rangle}\right)dx \\ &= \int_{\mathbb{R}^n} (-2 \pi i x_k) \left(f(x) e^{-2\pi i \langle x, y \rangle}\right)dx  =
     (-2 \pi ix_k f (x))\,\,\widehat{  } \,\, (y).\qedhere
\end{align}
\end{proof}

\begin{corolario}\label{coro:deri2}
Sea 
   $f \in \mathscr{L}^1(\mathbb{R}^n)$, y $n \in \mathbb{N}$. Supongamos que se cumple la condición 
   \begin{equation}\int_{\mathbb{R}^n}|x_{k}^{n}f(x)| \, dx < \infty,
   \end{equation}
   donde $x_k$ denota la coordenada k-ésima de $x$.
   Entonces, $\widehat{f}$ es $n$ veces derivable con respecto a $y_k$, y se verifica que
   \begin{equation}\label{eq:ecuacion}
    \frac{\partial \widehat{f}}{\partial y_k}(y) = ((-2 \pi ix_k)^nf(x))\,\,\widehat{  } \,\, (y) \quad \forall  y \in \mathbb{R}^n.
   \end{equation}
    Equivalentemente, para $\alpha = ne_k $,
   \begin{equation}
     D_{\alpha}\widehat{f}(y) = ((-2 \pi i)^{|\alpha|}x^{\alpha} f (x))\,\,\widehat{  } \,\, (y) \quad \forall  y \in \mathbb{R}^n.
   \end{equation}
\end{corolario}





\begin{proof}
Para cada $j \in \{1, \ldots ,n \}$, tenemos que:
\begin{equation}
    \int_{\mathbb{R}^n}|x_{k}^{j}f(x)| \, dx = \int_{|x_k|\leq 1}|x_{k}^{j}f(x)|dx+\int_{|x_k|>1}|x_{k}^{j}f(x)| \, dx\leq ||f||_1 +  \int_{\mathbb{R}^n }|x_{k}^{n}f(x)| \, dx < \infty.
\end{equation}

\noindent Luego se obtiene que $x_{k}^jf \in \mathscr{L}^1(\mathbb{R}^n)$. Realizando ahora una aplicación inductiva del teorema anterior, se demuestra el resultado buscado.
\end{proof}

\begin{teorema}\label{teo:der_trans}
   Sea $f \in \mathscr{L}^1(\mathbb{R}^n)$ y $\alpha = (\alpha_1, \ldots \alpha_n) \in (\mathbb{N} \cup \{0\})^n$. Supongamos que se cumple la condición 
   \begin{equation}\int_{\mathbb{R}^n}|x^{\beta}f(x)| \, dx < \infty  \quad \forall \beta \in A_{\alpha}.
   \end{equation}
   
\noindent Entonces, se tiene que
  
   \begin{equation}
     D_{\alpha}\widehat{f}(y) = ((-2 \pi i)^{|\alpha|}x^{\alpha} f (x))\,\,\widehat{  } \,\, (y) \quad \forall y \in \mathbb{R}^n.
   \end{equation}
\end{teorema} 


\begin{proof}
     Para un $\alpha \in ( \mathbb{N} \cup \{0\})^n$ cualquiera, basta con usar el Teorema \ref{teo:der_trans} y realizar inducción sobre $|\alpha|$.
\end{proof}



\section{Ejemplos}
\begin{definicion}\label{gauss}
    Definimos $G : \mathbb{R}^n \rightarrow \mathbb{R}$ función de $  \mathscr{L}^1(\mathbb{R}^n)$ dada por
    \begin{equation}
        G(x) = e^{- \pi ||x||^2} =  e^{- \pi\sum\limits_{k=1}^{n} x_k^2} \quad \forall x \in \mathbb{R}^n.
    \end{equation}
\end{definicion}


\begin{observacion}\label{obs:desc}
    Dado $x = (x_1,x_2,\ldots, x_n) \in \mathbb{R}^n$ se tiene que 
    \begin{equation}\label{eq:gauss}
        G(x_1,x_2,\ldots, x_n) = e^{- \pi x_1^2} \cdot e^{- \pi x_2^2} \cdots e^{- \pi x_n^2} = G(x_1)\cdot G(x_2) \cdots G(x_n)
    \end{equation}
\end{observacion}



\begin{lema}\label{lem:leam}
$G$ cumple
\begin{equation}
    \int_{\mathbb{R}^n} G(x)\, \, dx = 1 .
\end{equation}
\end{lema}
\begin{proof}
Por un lado sabemos que
\begin{equation}
   \widehat{G}(0,0, \ldots, 0) = \int_{\mathbb{R}^n} G(x)\, \, dx   .
\end{equation}

\noindent Por otro, atendiendo a la expresión \eqref{eq:gauss} y a la Proposición (\ref{prop:div}), se tiene que, 
\begin{equation}
    \widehat{G}(0, \ldots, 0) = \widehat{G}(0) \cdots \widehat{G}(0),
\end{equation}

\noindent luego podemos reducirnos al caso $n=1$.

\noindent Para ello, haremos un cálculo previo que nos servirá en nuestro objetivo.
\begin{equation}
    \int_{\mathbb{R}^2}e^{- \pi (x^2+y^2)} \, d(x,y) = \int_{-\pi}^{\pi}\int_{0}^{\infty}\rho e^{- \pi \rho^2} \, d \rho d\theta = \int_{-\pi}^{\pi} \frac{1}{2 \pi}d\theta = 1.
\end{equation}
Finalmente, 
\begin{equation}
\int_{\mathbb{R}^2}e^{- \pi (x^2+y^2)} \, d(x,y) = \int_{-\infty}^{\infty}\int_{-\infty}^{\infty} e^{- \pi x^2}e^{- \pi y^2} \, dy dx = \int_{-\infty}^{\infty}e^{- \pi x^2}  \, dx\int_{-\infty}^{\infty} e^{- \pi y^2} \, dy = \left( \int_{-\infty}^{\infty} G(t) \, dt \right)^2.
\end{equation}

\noindent De donde finalmente se deduce que para $G: \mathbb{R}  \rightarrow \mathbb{R}$
\begin{equation}
    \left( \int_{-\infty}^{\infty} G(t) \, dt \right)^2  = 1,
\end{equation}
y por ende 
\begin{equation}
    \widehat{G}(0) = \left( \int_{-\infty}^{\infty} G(t) \, dt \right)  = 1,
\end{equation}

\noindent concluyendo lo que queríamos.




\end{proof}



\begin{proposicion}
    $G$ cumple la identidad 
    \begin{equation}
        \widehat{G} = G.
    \end{equation}
\end{proposicion}




\begin{proof}
Por un lado, atendiendo a la expresión \eqref{eq:gauss} y a la Proposición (\ref{prop:div}), se tiene que, para $y = (y_1,y_2, \ldots, y_n)$,

\begin{equation}
    \widehat{G}(y_1,y_2, \ldots, y_n) = \widehat{G}(y_1) \cdots \widehat{G}(y_n). 
\end{equation}
\noindent Luego podemos reducirnos al caso n = 1.
\noindent Como 
\begin{equation}
    G'(x) = -2 \pi x G(x) \quad \forall x \in \mathbb{R},
\end{equation}
y además 
\begin{equation}
    \int_{\mathbb{R}}|G'(x)|dx = 2\int_0^\infty 2\pi x e^{-\pi x^2} \, dx = 2,
\end{equation}
\noindent el Teorema (\ref{teo:der_trans}) garantiza que $\widehat{G}$ es derivable en $\mathbb{R}$. También, 
\begin{equation}
    \widehat{G}'(y) = (-2 \pi i x G) \, \, \widehat{}\, \, (y) = i\widehat{G'}(y) = -2\pi y\widehat{G}(y) \quad \forall y \in \mathbb{R}.
\end{equation}
Luego $\widehat{G}$ es solución de la ecuación diferencial 
\begin{equation}
    f'(y) = -2\pi y f(y).
\end{equation}y, por tanto, 
\begin{equation}
    \widehat{G}(y) = ce^{-\pi y^2} = c G(y) \quad \forall y \in \mathbb{R},
\end{equation}
y usando el Lema (\ref{lem:leam})
\begin{equation}
    c = \widehat{G}(0) = \int_{\mathbb{R}}G(x) dx = 1.
\end{equation}
Obtenemos finalmente $\widehat{G} = G$.


\noindent Volvemos ahora al caso general. Sea  $y = (y_1,y_2,  \ldots, y_n)$,  se verifica

\begin{equation*}
        \widehat{G}(y) =  \prod_{j=1}^{n}\widehat{f}(y_j) =  \prod_{j=1}^{n}\widehat{f}(y_j)  = G(y). \qedhere
\end{equation*} 
\end{proof}

\begin{observacion}
  El proceso de filtrado de imágenes se emplea en el ámbito del procesamiento de imágenes con el objetivo de mejorar su calidad o de extraer información pertinente de las mismas. Consiste en aplicar diferentes operaciones matemáticas a los píxeles de una imagen con el fin de resaltar ciertas características, eliminar ruido o suavizar detalles no deseados. Este  filtro recibe el nombre de núcleo o (kernel). Dentro del filtrado de imágenes, el núcleo gaussiano, que aparecerá con frecuencia en los siguientes ejemplos, es usado frecuentemente.
    Este núcleo nace de la discretización de la función gaussiana 
    \begin{equation}
        G_{\sigma}(x,y) = \frac{1}{2\pi \sigma^2}e^{-\frac{x^2+y^2}{2 \sigma^2}} \quad \sigma \in \mathbb{R}^+.
    \end{equation}
    La distribución es extensamente conocida, y decae a medida que nos alejamos de su centro. Por tanto, como filtro discreto, le da más peso a los píxeles centrales y menos a los vecinos, ``difuminando'' los píxeles de la imagen. Por esto, es denominado un filtro de paso bajo o suavizado. El parámetro $\sigma$ controla la cantidad de suavizado.
    Nótese que 
    \begin{equation}\label{flip}
        G_{\sigma}(x) = \frac{1}{2\pi \sigma^2}H_{\frac{1}{\sqrt{\pi 2 \sigma^2}}}G(x) \quad \forall x \in \mathbb{R}^2.
    \end{equation}   
    De~\eqref{flip} deducimos las siguientes características  que simplifican el trabajo con la función gaussiana y contribuyen a su amplia utilización. 
    \begin{itemize}
        \item La Transformada de Fourier de una función Gaussiana es otra función Gaussiana
         \begin{equation}
        \widehat{G_{\sigma}}(y) = \frac{1}{2\pi \sigma^2} (2\pi \sigma^2)\widehat{G}((2\pi \sigma^2 y)) =G(\sqrt{2\pi \sigma^2}y) = G_{\sigma'}(y),
    \end{equation}
        donde $\sigma' = \frac{1}{2 \pi \sigma}$.
        \item Es separable, es decir se puede dividir en el producto de dos Gaussianas 1D.
    \end{itemize}
\end{observacion}
\noindent Cuando nos refiramos a la derivada del núcleo gaussiano, será una discretización de la derivada de la función gaussiana salvo constante.

\section{Teorema de Inversión}
A continuación abordamos el Teorema de Inversión que nos permitirá recuperar la función a partir de su Transformada de Fourier bajo ciertas hipótesis. 

\begin{teorema}[de Inversión]\label{teo:inv} Sea $f \in \mathscr{L}^1(\mathbb{R}^n)$ y supongamos que $\widehat{f}    \in \mathscr{L}^1(\mathbb{R}^n)$. Entonces
\begin{equation}\label{eq:teoinv}
    f(x) = \int_{\mathbb{R}^n}\widehat{f}(y)e^{2\pi i \langle x, y \rangle} \, dy \quad \text{para casi todo} \, x \in \mathbb{R}^n.
\end{equation}
   
\end{teorema}

\begin{proof}
Sea la función G definida en (\ref{gauss}). Dado $x \in \mathbb{R}^n$, probaremos primero que 
\begin{equation}\label{eq:primero}
   \underset{\substack{m \rightarrow \infty}}{\lim}\int_{\mathbb{R}^n}\widehat{f}(y)e^{2\pi i \langle x, y \rangle}G\left(\frac{1}{m}y\right) \, dy = \int_{\mathbb{R}^n}\widehat{f}(y)e^{2\pi i \langle x, y \rangle} \, dy .
\end{equation}
Para ello, observamos los siguientes puntos:
\begin{itemize}
    \item $\underset{\substack{m \rightarrow \infty}}{\lim}\, \widehat{f}(y)e^{2\pi i \langle x, y \rangle}G\left(\frac{1}{m}y\right)  = \widehat{f}(y)e^{2\pi i \langle x, y \rangle}$,
    \item $|\widehat{f}(y)e^{2\pi i \langle x, y \rangle}G\left(\frac{1}{m}y\right)| \leq |\widehat{f}(y)| \quad \forall y \in \mathbb{R}^n, \forall m \in \mathbb{N}$,
    \item $|\widehat{f}|$ es por hipótesis integrable en $\mathbb{R}^n$.
\end{itemize}

\noindent Aplicando el Teorema de la convergencia dominada se concluye~\eqref{eq:primero}.

\noindent Por otro lado, usando el Lema \ref {lem:leam} y el Teorema \ref{teo:cambio}, se tiene que
\begin{align}
& \int_{\mathbb{R}^n}\widehat{f}(y)e^{2\pi i \langle x, y \rangle}G\left(\frac{1}{m}y\right) \, dy - f(x) = \\
& = \int_{\mathbb{R}^n}\widehat{f}(y)e^{2\pi i \langle x, y \rangle}G\left(\frac{1}{m}y\right) \, dy - \int_{\mathbb{R}^n}G(y)f(x)\, dy \\
& = \int_{\mathbb{R}^n}\widehat{ \tau_{-x}f}(y)G\left(\frac{1}{m}y\right) \, dy - \int_{\mathbb{R}^n}G(y)f(x)\, dy \\
& = \int_{\mathbb{R}^n} \tau_{-x}f(y)\widehat{G} \left(\frac{1}{m}y\right) \, dy - \int_{\mathbb{R}^n}G(y)f(x)\, dy \\
& = \int_{\mathbb{R}^n} \tau_{-x}f(y) m^mG \left(my\right) \, dy - \int_{\mathbb{R}^n}G(y)\, dy \\
& = \int_{\mathbb{R}^n}[ f(\frac{y}{m}+x)-f(x)] G \left(y\right)\, dy.
\end{align}
Luego
\begin{equation}
    \left|\int_{\mathbb{R}^n}\widehat{f}(y)e^{2\pi i \langle x, y \rangle}G\left(\frac{1}{m}y\right) \, dy - f(x)\right| \leq  \int_{\mathbb{R}^n}\left|f(\frac{y}{m}+x)-f(x) \right| G \left(y\right)\, dy.
\end{equation}
Integrando con respecto a $x \in \mathbb{R}^n$ se tiene
\begin{align}
\int_{\mathbb{R}^n}\left|\int_{\mathbb{R}^n}\widehat{f}(y)e^{2\pi i \langle x, y \rangle}G\left(\frac{1}{m}y\right) \, dy - f(x)\right| \, dx &\leq \int_{\mathbb{R}^n}\int_{\mathbb{R}^n}\left|f(\frac{y}{m}+x)-f(x) \right| G \left(y\right)\, dx dy \nonumber \\
&= \int_{\mathbb{R}^n} w_1f\left(\frac{1}{m}y\right)G(y)\, dy.
\end{align}
Recurriendo de nuevo al Teorema de la convergencia dominada, probaremos que 
\begin{equation}
    \lim_{m \to \infty}  \int_{\mathbb{R}^n} w_1f\left(\frac{1}{m}y\right)G(y)\, dy = 0.
\end{equation}
Observamos los siguientes puntos:
\begin{itemize}
    \item $\underset{\substack{m \rightarrow \infty}}{\lim}\, w_1f\left(\frac{1}{m}y\right)G(y)  = 0$,
    \item $w_1f\left(\frac{1}{m}y\right)G(y) \leq 2||f||_1G(y) \quad \forall y \in \mathbb{R}^n, \forall m \in \mathbb{N}$,
    \item La función $||f||_1G$ es integrable en $\mathbb{R}^n$.
\end{itemize}
Por lo que se verifican todas las hipótesis para poder aplicar de nuevo el Teorema de la convergencia dominada para obtener que
\begin{equation}
    \lim_{m \to \infty}\int_{\mathbb{R}^n}\left|\int_{\mathbb{R}^n}\widehat{f}(y)e^{2\pi i \langle x, y \rangle}G\left(\frac{1}{m}y\right) \, dy - f(x)\right| \, dx = 0.
\end{equation}
Esto, junto con~\eqref{eq:primero}, nos permite concluir que 
\begin{equation*}
    \int_{\mathbb{R}^n}\widehat{f}(y)e^{2\pi i \langle x, y \rangle} \, dy =f(x) \quad c.p.d. \,\, \text{en}\,\, {\mathbb{R}^n}. \qedhere
\end{equation*}
    
\end{proof}


\begin{corolario}
    La fórmula que aparece en~\eqref{eq:teoinv}  se cumple para cualquier punto en el que $f$ es continua. 
\end{corolario}
\begin{proof}
    Dado $a \in \mathbb{R}$, supongamos que f es continua en $a$ y razonemos por reducción al absurdo.
    Si $F(a)\neq f(a)$, entonces $\exists \delta \in \mathbb{R^+}$ tal que $F(x)\neq f(x) \, \forall x \in ]a-\delta, a+\delta[$, por ser $F$ y $f$ funciones continuas en el punto $a$.
    Como consecuencia, existiría un conjunto de medida no nula $]a-\delta, a+\delta[$ en el que no se verifica~\eqref{eq:teoinv} contrariando lo obtenido en el Teorema \ref{teo:inv}.
\end{proof}


\begin{observacion}
   El Teorema de Inversión afirma que 
   \begin{equation}
       f(x) = \widehat{\widehat{f}(}-x )\,\, c.p.d \,\, \text{en} \,\, \mathbb{R}^n.
   \end{equation}
Y como la función que aparece en la derecha es continua en $\mathbb{R}^n$, deducimos que $f$ no puede estar acotada en un entorno reducido del punto y ser discontinua en este. ¿?
\end{observacion}

\noindent Es interesante puntualizar que la condición de que la Transformada de Fourier sea integrable, es bastante restrictiva y en general no será cierta.

\noindent Sin embargo en el ámbito discreto, como es de esperar, podremos recuperar cualquier señal a partir de la Trasformada usando la Inversa de la Transformada de Fourier Discreta (IDFT). 





\section{Teorema de Unicidad}
El Teorema de Unicidad sugiere que la Transformada de Fourier $\widehat{f}$ puede interpretarse como el código genético de la función $f$. En efecto, si dos funciones integrables en $\mathbb{R}^n$ tienen la misma Transformada de Fourier, esto implicará que son idénticas salvo en un conjunto de medida nula.
Por tanto, cabe pensar que la Transformada de Fourier, en cierto sentido, sintetiza la función original.





\begin{teorema}(Teorema de Unicidad).
Sean $f,g \in \mathscr{L}^1(\mathbb{R}^n)$ y supongamos que $\widehat{f}(y) = \widehat{g}(y)$ para todo $ y \in \mathbb{R}^n$. Entonces $f = g$ c.p.d.  en $\mathbb{R}^n$.
\end{teorema}

\begin{proof}
La prueba se deduce del Teorema de Inversión.

\noindent Si $f,g \in \mathscr{L}^1(\mathbb{R}^n)$, entonces $f-g \in \mathscr{L}^1(\mathbb{R}^n)$ con $\widehat{f}(y) - \widehat{g}(y) = \widehat{f-g}(y) = 0 \quad \forall y \in \mathbb{R}^n$.
Por tanto, el Teorema de Inversión es aplicable a la función $f-g$ y deducimos que
\begin{equation}
    f(x)-g(x) = \int_{\mathbb{R}^n}(\widehat{f-g})(y)e^{2 \pi i\langle x, y \rangle} \, dy = 0. \quad \text{c.p.d. en } \mathbb{R}^n ,
\end{equation}
de donde deducimos que $f=g$ c.p.d en $\mathbb{R}^n$.
\end{proof}

\section{Clase de Schwartz}
Tal como se ha señalado previamente, existen referencias, tales como \cite{LibroSchwartz} donde se desarrolla la Transformada de Fourier en una clase más restringida que $\mathscr{L}^1(\mathbb{R}^n)$, el espacio de funciones de Schwartz $\mathscr{S}(\mathbb{R}^n)$,
que resulta ser un entorno particularmente adecuado para esta. Es por ello por lo que estudiaremos algunas características de este espacio, así como la razón de que el comportamiento de la Transformada sea especialmente bueno en él.

\begin{definicion}\label{def111}

\noindent El espacio de Schwartz, denotado por \( \mathscr{S}(\mathbb{R}^n) \), se define como el conjunto de todas las funciones \( f: \mathbb{R}^n \rightarrow \mathbb{C} \) que son infinitamente diferenciables, de modo que para cada par de multi-índices $\alpha$ y $\beta \in (\mathbb{N}\cup 0)^n$, se tiene que
\[
\rho_{\beta, \alpha}(f) := \sup_{x \in \mathbb{R}^n} |x^\beta D_\alpha f(x)| < \infty.
\]
\end{definicion}
\begin{observacion}
    Las funciones de \( \mathscr{S}(\mathbb{R}^n) \) son funciones infinitamente diferenciables cuyas sucesivas derivadas decrecen más rápido en infinito que cualquier polinomio.
\end{observacion}

\noindent Es posible proporcionar otra definición utilizando un único multi-índice.

\begin{definicion}
Sea \(f \in C^{\infty}(\mathbb{R}^n)\). Se dice que \(f\) pertenece al espacio de Schwartz si es indefinidamente derivable y para cada \(m \in \mathbb{N}\) y cada multi-índice \(\alpha \in (\mathbb{N}\cup \{0\})^n\) se cumple que
\[
\rho_{m, \alpha}(f) := \sup_{x \in \mathbb{R}^n} (1+\|x\|^2)^m |D_\alpha f(x)|  < \infty.
\]
\end{definicion}
\noindent Es claro que ambas definiciones son equivalentes. Se ha usado a propósito la misma notación para los dos supremos.

\begin{proposicion}\label{profdif1}
    Si $f \in \mathscr{S}(\mathbb{R}^n)$. Entonces $f \in \mathscr{L}^1(\mathbb{R}^n)$.
\end{proposicion}

\begin{proof}
Tomamos $\alpha = 0$. Como $f \in \mathscr{S}(\mathbb{R}^n)$, se tiene que 
\begin{equation}
    |f(x)|(1+\|x\|^2)^n \leq \rho_{n,\alpha},
\end{equation}
Usando coordenadas polares se tiene que
\begin{equation*}
\begin{aligned}
   \int_{\mathbb{R}^n}|f(x)| \, dx &= \int_{\mathbb{R}^n}\frac{\rho_{n,\alpha}}{(1+\|x\|^2)^n} \, dx = \rho_{n,\alpha}w_{n-1}\int_{0}^{\infty} \frac{r^{n-1}}{(1+r^2)^n} \, dr < \infty,
   \end{aligned}
\end{equation*} 
\noindent donde $w_{n-1}$ denota el área de $\mathbb{S}^{n-1}$.
\end{proof}

\begin{proposicion}\label{prodif2}
    Si $f,g \in \mathscr{S}(\mathbb{R}^n)$, entonces se tiene que 
    \begin{itemize}
        \item $f+g \in \mathscr{S}(\mathbb{R}^n)$,
        \item $\gamma f \in \mathscr{S}(\mathbb{R}^n) \quad \forall \gamma \in \mathbb{C}$.
    \end{itemize}
\end{proposicion}

\begin{proof}
Probaremos a continuación los dos apartados.
    \begin{itemize}
        \item Veamos que $f+g \in \mathscr{S}(\mathbb{R}^n)$. Para cada par de multi-índices $\alpha, \beta \in (\mathbb{N} \cup {0})^n$, se tiene que 
        \begin{equation}
            D_{\alpha}(f+g) =   D_{\alpha}(f)+ D_{\alpha}(g).
        \end{equation}
        Luego
       \begin{align}
        |x^{\beta}D_{\alpha}(f+g)| 
        &= |x^{\beta}||D_{\alpha}(f)+ D_{\alpha}(g)| \\
        &\leq |x^{\beta}||D_{\beta}(f)| + |x^{\beta}||D_{\beta}(g)| \\
        &\leq \rho_{\beta,\alpha}(f)+\rho_{\beta,\alpha}(g).
        \end{align}
        Y, por tanto, 
        \begin{equation}
            \rho_{\beta,\alpha}(f+g) \leq \rho_{\beta,\alpha}(f)+\rho_{\beta,\alpha}(g) < \infty.
        \end{equation}
        Concluimos que $f+g \in \mathscr{S}(\mathbb{R}^n)$.

        \item Sea $\gamma \in \mathbb{C}$. Veamos que $\gamma f \in \mathscr{S}(\mathbb{R}^n)$.  Para cada par de multi-índices $\alpha, \beta \in (\mathbb{N} \cup {0})^n$, se tiene que 
        \begin{equation}
            D_{\alpha}(\gamma f) =  \gamma D_{\alpha}(f).
        \end{equation}
        Luego
       \begin{equation}
        |x^{\beta}D_{\alpha}(\gamma f)| 
        = |x^{\beta}\gamma D_{\alpha}(f)|
        = |\gamma||x^{\beta}D_{\alpha}(f)| 
        \leq |\gamma|\rho_{\beta,\alpha}(f).
        \end{equation}
        Y, por tanto, 
        \begin{equation}
        \rho_{\beta,\alpha}(\gamma f) \leq |\gamma|\rho_{\beta,\alpha}(f) < \infty,
        \end{equation}
        obteniendo que $\gamma f \in \mathscr{S}(\mathbb{R}^n)$. \qedhere


        
        \end{itemize}        
\end{proof}
    



\begin{observacion}
\noindent Las proposiciones \ref{profdif1} y \ref{prodif2} muestran que el espacio $\mathscr{S}(\mathbb{R}^n)$ es un subespacio vectorial de  $\mathscr{L}^1(\mathbb{R}^n)$.
Consecuentemente, la teoría desarrollada para $\mathscr{L}^1(\mathbb{R}^n)$ se aplica directamente a este espacio.
\end{observacion}

\begin{proposicion}\label{propopes}
    Si $f \in \mathscr{S}(\mathbb{R}^n)$. Entonces $\widehat{f} \in \mathscr{S}(\mathbb{R}^n)$.
\end{proposicion}

\begin{proof}
    Probaremos primero que $\widehat{f}$ es infinitamente derivable. Como puede intuir el lector, pretendemos usar el Teorema \ref{teo:der_trans}. Para ello tenemos que demostrar que si $f \in \mathscr{S}(\mathbb{R}^n)$, ser verifican las hipótesis de este teorema.
    
\noindent  Sea $\alpha = (\alpha, \ldots \alpha_n)$, atendiendo a la equivalencia de las dos definiciones de $\mathscr{S}(\mathbb{R}^n)$ proporcionadas se tiene que 
\begin{equation}
\begin{aligned}
   \int_{\mathbb{R}^n}|x^{\alpha}f(x)| \, dx 
   &< \infty.
\end{aligned}
\end{equation}
\noindent Por tanto, se tiene que $\widehat{f}$ es infinitamente derivable y además
      \begin{equation}
     D_{\alpha}\widehat{f}(y) = ((-2 \pi i)^{|\alpha|}x^{\alpha} f (x))\,\,\widehat{  } \,\, (y) \quad \forall y \in \mathbb{R}^n.
   \end{equation}
Definimos la función 
\begin{equation}
g(x) = (-2 \pi i)^{|\alpha|}x^{\alpha} f (x) \quad \forall x \in \mathbb{R}^n.
\end{equation}

\noindent Dado $\beta \in  (\mathbb{N} \cup {0})^n $,
    \begin{equation}
     y^{\beta}D_{\alpha}\widehat{f}(y) =  y^{\beta}\widehat{g(x)}(y) = \widehat{D_{\beta}g(x)}(y) \frac{1}{(2 \pi i)^{|\beta|}},
    \end{equation}
de donde se tiene que $y \mapsto y^{\beta}D_{\alpha}\widehat{f}$ está acotada $\forall \beta \in  (\mathbb{N} \cup {0})^n $, ya que es la Transformada de Fourier de una función integrable en $\mathbb{R}^n$. 
Como consecuencia deducimos que  $\widehat{f} \in \mathscr{S}(\mathbb{R}^n)$. 
\end{proof} 





\begin{observacion}
Trabajando el concepto de la Transformada
de Fourier en $\mathscr{S}(\mathbb{R}^n)$, tenemos garantizado
de antemano el buen planteamiento de su construcción inversa, ya que para poder aplicar el Teorema de Inversión a una función $f$ necesitábamos que $\widehat{f}$ fuera integrable, una condición que gracias a la proposición anterior se cumple para  $f \in \mathscr{S}(\mathbb{R}^n)$.
\end{observacion}

\begin{proposicion}
    El operador
    \begin{equation}
        \widehat{} :  \mathscr{S}^1(\mathbb{R}^n) \rightarrow \mathscr{S}^1(\mathbb{R}^n)
    \end{equation}
define una biyección.
\end{proposicion}
\begin{proof}
\noindent Sabemos por la Proposición \ref{propopes} que  $\widehat{f} \in \mathscr{S}(\mathbb{R}^n)$ para cada $f \in \mathscr{S}(\mathbb{R}^n)$. Además,
\begin{itemize}
    \item Es inyectivo como consecuencia del Teorema de Unicidad. Esto es: si  $\widehat{f}=0$, se tiene que $f=0$ c.p.d.
    \item Es sobreyectivo por el Teorema de Inversión.  Dada una función $g \in \mathscr{S}(\mathbb{R}^n)$, definimos la función $f = \widetilde{\widehat{g}}$, de tal manera que por el Teorema de Inversión se tiene justamente que $\widehat{f}= g$.
\end{itemize}

\end{proof}






\noindent La familia numerable de seminormas 
\[\sup_{|\alpha|\leq N} \sup_{x \in \mathbb{R}^n}(1+||x||_2^2)^k |D_\alpha f(x)| \quad (f\in \mathscr{S}(\mathbb{R}^d),k,N \in \mathbb{N}),\]
da lugar a una topología en $S(\mathbb{R}^n)$, localmente convexa y metrizable, generada por intersecciones finitas de bolas definidas con respecto a tales seminormas. El dual topológico de $S(\mathbb{R}^n)$, denotado $S'(\mathbb{R}^n)$, está formado por todos los funcionales lineales y continuos definidos sobre $S(\mathbb{R}^n)$, y sus elementos reciben el nombre de distribuciones temperadas. En este espacio se puede dar una definición de la Transformada de Fourier y probar que la Transformada
de Fourier de una distribución temperada es otra distribución temperada. También se puede probar que la Transformada de Fourier es un homeomorfismo lineal de
$S(\mathbb{R}^n)$ en $S(\mathbb{R}^n)$, y de $S'(\mathbb{R}^n)$ en $S'(\mathbb{R}^n)$. Esta aproximación, aunque escapa del alcance del presente análisis, se presenta como una línea prometedora para trabajo futuro.



%--------------------------------------------------------------------
% FIN DEL CAPÍTULO. 
%--------------------------------------------------------------------


% Información relevante para la elaboración del trabajo.
%% !TeX root = ../tfg.tex
% !TeX encoding = utf8

\chapter{Documentación}\label{ch:primer-capitulo}

\section{Introducción}
Este documento es una plantilla para la elaboración de un trabajo fin de Grado siguiendo los \href{https://grados.ugr.es/matematicas/pages/infoacademica/tfg/requisitosTFG}{requisitos} de la comisión de Grado en Matemáticas de la Universidad de Granada que, a fecha de junio de 2023, son las siguientes:

\begin{itemize}
  \item La  memoria  debe  realizarse  con  un  procesador  de  texto  científico,  preferiblemente (La)TeX.
  \item La portada  debe contener  el  logo  de  la UGR,  incluir  el  título del TFG, el nombre del estudiante y especificar el grado, la facultad y el curso actual.
  \item La contraportada contendrá además el nombre del tutor o tutores.
  \item La memoria debe necesariamente incluir:
    \begin{itemize}
      \item Declaración explícita firmada en la que se asume la originalidad del trabajo, entendida en el sentido de que no ha utilizado fuentes sin citarlas debidamente. Esta declaración se puede descargar en la web del Grado.
      \item un índice detallado de capítulos y secciones,
      \item un resumen amplio en inglés del trabajo realizado (se recomienda entre 800 y 1500 palabras),
      \item una introducción en la que se describan claramente los objetivos previstos inicialmente en la propuesta de TFG, indicando si han sido o no alcanzados, los antecedentes importantes para el desarrollo, los resultados obtenidos, en su caso y las principales fuentes consultadas,
      \item una bibliografía final que incluya todas las referencias utilizadas.
    \end{itemize}
  \item Se recomienda que la extensión de la memoria sea de unas 50 páginas, sin incluir posibles apéndices.
\end{itemize}

Para generar el pdf a partir de la plantilla basta compilar el fichero \texttt{tfg.tex}. Es conveniente leer los comentarios contenidos en dicho fichero pues ayudarán a entender mejor como funciona la plantilla. 

La estructura de la plantilla es la siguiente\footnote{Los nombres de las carpetas no se han acentuado para evitar problemas en sistemas con Windows}: 
\begin{itemize}
  \item Carpeta \textbf{preliminares}: contiene los siguientes archivos
    \begin{description}
      \item[\texttt{dedicatoria.tex}] Para la dedicatoria del trabajo (opcional)
      \item[\texttt{agradecimientos.tex}] Para los agradecimientos del trabajo (opcional)
      \item[\texttt{introduccion.tex}] Para la introducción (obligatorio)
      \item[\texttt{summary.tex}] Para el resumen en inglés (obligatorio)
      \item[\texttt{tablacontenidos.tex}] Genera de forma automática la tabla de contenidos, el índice de figuras y el índice de tablas. Si bien la tabla de contenidos es conveniente incluirla, el índice de figuras y tablas es opcional. Por defecto está desactivado. Para mostrar dichos índices hay que editar este fichero y quitar el comentario a \verb+\listoffigures+ o \verb+\listoftables+ según queramos uno de los índices o los dos. En este archivo también es posible habilitar la inclusión de un índice de listados de código (si estos han sido incluidos con el paquete \texttt{listings})
  \end{description}
  El resto de archivos de dicha carpeta no es necesario editarlos pues su contenido se generará automáticamente a partir de los metadatos que agreguemos en \texttt{tfg.tex}

  \item Carpeta \textbf{capitulos}: contiene los archivos de los capítulos del TFG. Añadir tantos archivos como sean necesarios. Este capítulo es \texttt{capitulo01.tex}.

  \item Carpeta \textbf{apendices}: Para los apéndices (opcional)
  \item Carpeta \textbf{img}: Para incluir los ficheros de imagen que se usarán en el documento.
    
  \item Fichero \texttt{library.bib}: Para incluir las referencias bibliográficas en formato \texttt{bibtex}. Es útil la herramienta \href{https://www.doi2bib.org/}{doi2bib} para generar de forma automática el código bibtex de una referencia a partir de su \textsc{doi}  así como la base de datos bibliográfica \href{https://mathscinet.ams.org}{MathSciNet}. Para que una referencia aparezca en el pdf no basta con incluirla en el fichero \texttt{library.bib}, es necesario además \emph{citarla} en el documento usando el comando \verb+\cite+. Si queremos mostrar todos las referencias incluidas en el fichero \texttt{library.bib} podemos usar \verb+\cite{*}+ aunque esta opción no es la más adecuada. Se aconseja que los elementos de la bibliografía estén citados al menos una vez en el documento (y de esa forma aparecerán de forma automática en la lista de referencias).

  \item Fichero \texttt{glosario.tex}: Para incluir un glosario en el trabajo (opcional). Si no queremos incluir un glosario deberemos borrar el comando \verb+% !TeX root = ../tfg.tex
% !TeX encoding = utf8

\chapter*{Glosario}
\addcontentsline{toc}{chapter}{Glosario} % Añade el glosario a la tabla de contenidos

La inclusión de un glosario es opcional.

Archivo: \texttt{glosario.tex}

\begin{description} 
  \item[$\mathbb{R}$] Conjunto de números reales.

  \item[$\mathbb{C}$] Conjunto de números complejos.

  \item[$\mathbb{Z}$] Conjunto de números enteros.
\end{description}
\endinput
+ del fichero \texttt{tfg.tex} y posteriormente borrar el fichero \texttt{glosario.tex}

   \item Fichero \texttt{tfg.tex}: El documento maestro del TFG que hay que compilar con \LaTeX\ para obtener el pdf. En dicho documento hay que cambiar la \emph{información del título del \textsc{tfg} y el autor así como los tutores}.
\end{itemize}



\section{Elementos del texto}

En esta sección presentaremos diferentes ejemplos de los elementos de texto básico. Conviene consultar el contenido de \texttt{capitulos/capitulo01.tex} para ver cómo se han incluido.

\subsection{Listas}
En \LaTeX\ tenemos disponibles los siguientes tipos de listas:

Listas enumeradas:
\begin{enumerate}
  \item item 1
  \item item 2
  \item item 3
\end{enumerate}

Listas no enumeradas
\begin{itemize}
  \item item 1
  \item item 2
  \item item 3
  \end{itemize}

Listas descriptivas
\begin{description}
  \item[termino1] descripción 1
  \item[termino2] descripción 2
\end{description}
  
\subsection{Tablas y figuras}

En la \autoref{tb:ejemplo-tabla} o la \autoref{fig:logo-ugr} podemos ver\ldots

\begin{table}[htpb]
  \centering
  \begin{tabular}{ccc} \toprule
    \multicolumn{2}{c}{Agrupados} \\ \cmidrule(r){1-2}
    cabecera & cabecera & cabecera          \\ \midrule
    elemento & elemento & elemento          \\ 
    elemento & elemento & elemento          \\ 
    elemento & elemento & elemento          \\ \bottomrule
  \end{tabular}
  \caption{Ejemplo de tabla}
  \label{tb:ejemplo-tabla}
\end{table}

\begin{figure}[htpb]
  \centering
  \includegraphics[width=0.5\textwidth]{logo-ugr}
  \caption{Logotipo de la Universidad de Granada}
  \label{fig:logo-ugr}
\end{figure}

\section{Entornos matemáticos}\label{sec:entornos-matematicos}

La plantilla tiene definidos varios entornos matemáticos cuyo nombre es el mismo omitiendo los acentos. Así, para incluir una \emph{proposición} usaríamos:

\begin{verbatim}
\begin{proposicion}
texto de la proposición
\end{proposicion} 
\end{verbatim}

Ver el código fuente del archivo \texttt{documentacion.tex} en la carpeta \texttt{capitulos} para el resto de ejemplos.

\begin{teorema}\label{thm:teorema}
Esto es un ejemplo de teorema.
\end{teorema}

\begin{proposicion}
Ejemplo de proposición
\end{proposicion}

\begin{lema}
Ejemplo de lema
\end{lema}

\begin{corolario}
Ejemplo de corolario
\end{corolario}

\begin{definicion}
Ejemplo de definición
\end{definicion}

\begin{observacion}
Ejemplo de observación
\end{observacion}

Adicionalmente está definido el entorno \texttt{teorema*} que permite incluir un teorema sin numeración:

\begin{teorema*}[Fórmula de Gauß-Bonnet]
  Sea $S$ una superficie compacta y $K$ su curvatura de Gauß. Entonces
\begin{equation}
  \int_S K = 2\pi\chi(S)
\end{equation}
donde $\chi(S)$ es la característica de Euler de $S$.
\end{teorema*}

Las fórmulas matemáticas se escriben entre símbolos de dólar \$ si van en línea con el texto o bien usando el entorno%
\footnote{
  También es posible delimitar una ecuación mediante los comandos \texttt{$\backslash$[} y \texttt{$\backslash$]} pero éstas nunca llevarán numeración aunque añadamos una etiqueta y las referenciemos (ver \autoref{sec:referencias}).
} 
\texttt{equation} cuando queremos que se impriman centradas en una línea propia, como el siguiente ejemplo
\begin{equation}\label{eq:identidad-pitagorica}
  \cos^2 x + \sin^2 x = 1.
\end{equation}


Gracias al paquete \texttt{mathtools}, las ecuaciones escritas dentro del entorno \texttt{equation} llevarán numeración de forma automática si son referenciadas  en cualquier parte del documento (por ejemplo la identidad Pitagórica~\eqref{eq:identidad-pitagorica}, ver el código de los dos anteriores ejemplos y la \autoref{sec:referencias} para más información sobre referencias cruzadas en el documento).




\section{Referencias a elementos del texto}\label{sec:referencias}

Para las referencias a los elementos del texto (secciones, capítulos, teoremas,\ldots) se procede de la siguiente manera:
\begin{itemize}
  \item Se \emph{marca} el elemento (justo después del mismo si se trata de un capítulo o sección o en el interior del \emph{entorno} en otro caso), mediante el comando \verb+\label{+\emph{etiqueta}\verb+}+, donde \emph{etiqueta} debe ser un identificador único.
  \item Para crear una referencia al elemento en cualquier otra parte del texto se usa el comando \verb+\ref{+\emph{etiqueta}\verb+}+ (únicamente imprime la numeración asociada a dicho elemento, por ejemplo \ref{ch:primer-capitulo} o \ref{thm:teorema}) o bien \verb+\autoref{+\emph{etiqueta}\verb+}+ (imprime la numeración del elemento así como un texto previo indicando su tipo, por ejemplo \autoref{ch:primer-capitulo} o \autoref{thm:teorema})
\end{itemize}




\section{Bibliografía e índice}

Esto es un ejemplo de texto en un capítulo. Incluye varias citas tanto a libros~\cite{Aigner2018}, artículos de investigación~\cite{Euler1985}, recursos online~\cite{EulerWiki} (páginas web), tesis~\cite{CitekeyPhdthesis}, trabajo fin de máster~\cite{CitekeyMastersthesis}, trabajo fin de grado~\cite{CiteKeyBachelorsthesis} así como artículos sin publicar (preprints) \cite{castroinfantes2022conjugate} (en estos últimos usar el campo \texttt{note} para añadir la información relevante). 

Ver el fichero \texttt{library.bib} para las distintas plantillas. Cada nueva referencia debe añadirse en dicho fichero siguiendo el estilo del código \texttt{bibtex} según el tipo de referencia (página web, tesis, trabajo fin de grado o máster, artículo de investigación, libro,\ldots). Alternativamente se puede usar la web \href{https://zbib.org}{https://zbib.org} para generar automáticamente el código \texttt{bibtex}.


\endinput

%\chapter[Consideraciones elaboración TFG]{Consideraciones generales para la elaboración de un trabajo fin de grado}

\section{Normativa de la comisión del Grado en Matemáticas}

El \textsc{tfg} lo rigen dos normativas: 
\begin{itemize}
  \item una a nivel general de la UGR (\href{https://secretariageneral.ugr.es/sites/webugr/secretariageneral/public/inline-files/BOUGR/187/PLANTILLA%20CABECERASDoc2.pdf}{Reglamento del Trabajo o Proyecto fin de Grado de la Universidad de Granada}\footnote{\url{https://secretariageneral.ugr.es/sites/webugr/secretariageneral/public/inline-files/BOUGR/187/PLANTILLA\%20CABECERASDoc2.pdf}}) y 
  \item otra complementaria a nivel de la Facultad de Ciencias  (\href{https://fciencias.ugr.es/images/stories/documentos/reglamentos/reglamentoTfgCiencias23.pdf}{Reglamento del trabajo fin de grado en la Facultad de Ciencias de la Universidad de Granada}\footnote{\url{https://fciencias.ugr.es/images/stories/documentos/reglamentos/reglamentoTfgCiencias23.pdf}}). 
\end{itemize} 
Además, la comisión del Grado de Matemáticas impone unos \href{https://grados.ugr.es/matematicas/pages/infoacademica/tfg/requisitosTFG}{Requisitos de la memoria}\footnote{\url{https://grados.ugr.es/matematicas/pages/infoacademica/tfg/requisitosTFG}}.

El \textsc{tfg} hay que elaborarlo preferiblemente en LaTeX y puede usar la plantilla disponible en \href{https://github.com/latex-mat-ugr/Plantilla-TFG/archive/master.zip}{Plantilla \textsc{tfg} grado en matemáticas formato .tex}\footnote{\url{https://github.com/latex-mat-ugr/Plantilla-TFG/archive/master.zip}}.

Toda la información anterior puede encontrarse en la \href{https://grados.ugr.es/matematicas/pages/infoacademica/trabajofingrado}{web del Grado en Matemáticas}\footnote{\url{https://grados.ugr.es/matematicas/pages/infoacademica/trabajofingrado}}.

Es conveniente tener presente la documentación anterior para la elaboración del \textsc{tfg}. En especial en lo relativo a las fechas de depósito del \textsc{tfg} para su defensa.

A continuación destaco algunos aspectos importantes de la misma:
\begin{itemize}
\item El plagio, entendido como la presentación de un trabajo u obra hecho por otra persona como propio o la copia de textos sin citar su procedencia y dándolos como de elaboración propia, conllevará automáticamente la calificación numérica de cero. Esta consecuencia debe entenderse sin perjuicio de las responsabilidades disciplinarias en las que pudieran incurrir los estudiantes que plagien.
  \item Las memorias entregadas por parte de los estudiantes tendrán que ir firmadas sobre una declaración explícita en la que se asume la originalidad del trabajo, entendida en el sentido de que no ha utilizado fuentes sin citarlas debidamente.
\end{itemize}


% Los criterios de evaluación utilizados permitirán evaluar el grado de adquisición de las competencias que tiene establecidas el TFG en el VERIFICA de la titulación. Además, entre otros aspectos, se tendrá en consideración:

% \begin{itemize}
%   \item Redacción y ortografía tanto en la memoria del TFG como en los medios usados para la defensa del mismo (diapositivas, etc.).
%   \item Adecuación al formato de memoria indicado. Se proporcionará una plantilla de memoria de TFG a tal fin.
%   \item Adecuación temporal a los cronogramas de trabajo según los plazos de entrega marcados por el tutor/es.
%   \item Nivel de profundidad en los contenidos expuestos.
%   \item Dominio del tema e iniciativa del alumno.
%   \item Claridad de la exposición y adecuación al tiempo de exposición establecido.
%   \item Capacidad de análisis y síntesis.
%   \item Discusión con la Comisión Evaluadora en el turno de preguntas.
% \end{itemize}

% Se entregará una copia escrita, a doble cara, de la memoria del TFG para su evaluación, por parte de la Comisión Evaluadora, con una antelación de una semana (7 días naturales) antes de la fecha de defensa pública del TFG. Además, se entregará una versión en formato "pdf" de dicha memoria que quedará en la base de datos de todos los TFG y que custodiará la CTFGOO.




\section{Formato de la memoria}
La memoria se presentará usando un editor de textos científico, preferiblemente \LaTeX, e incluir los siguientes apartados:
\begin{enumerate}
  \item \emph{Resumen en inglés}: Deberá estar escrito completamente en inglés y tener una longitud recomendada entre 800 y 1500 palabras. 
  \item \emph{Introducción}. Deberá:
    \begin{itemize}
      \item Indicar los \emph{Objetivos del trabajo}: deberán aparecer con claridad los objetivos inicialmente previstos en la propuesta de \textsc{tfg} y los finalmente alcanzados con indicación de dificultades, cambios y mejoras respecto a la propuesta inicial. Si procede, es conveniente apuntar de manera precisa las interdependencias entre los distintos objetivos y conectarlos con los diferentes apartados de la memoria. Se pueden destacar aquí los aspectos formativos previos más utilizados. 
    \item Contextualizar el trabajo explicando antecedentes importantes para el desarrollo realizado y efectuando, en su caso, un estudio de los progresos recientes.
    \item Describir el problema abordado, de forma que el lector tenga desde este momento una idea clara de la cuestión a resolver o del producto a desarrollar y una visión general de la solución alcanzada.
    \item Indicar los resultados obtenidos.
    \item Citar las principales fuentes consultadas.
    \end{itemize}

  \item \emph{Desarrollo del trabajo}: El trabajo se estructurará en partes o capítulos según convengan, con la posibilidad de incluir apéndices. Se recomienda que la extensión de esta parte (sin incluir los apéndices) sea de unas 50 páginas. 

  \item \emph{Conclusiones y vías futuras}: Las conclusiones deberán incluir todas aquellas de tipo profesional y académico. Si hubiese posibles vías claras de desarrollo posterior sería interesante destacarlas aquí, poniéndolas en valor en el contexto inicial del trabajo.

  \item \emph{Bibliografía final}: Se incluirán tanto las fuentes primarias como todas aquellas cuyo peso haya sido menor en la realización del trabajo. Se recomienda un breve comentario de las referencias, ya sea individualizado, por grupos de referencias o global. En caso de incluir \textsc{url}s de páginas web deberán ir acompañadas de título, autor y fecha de último acceso, entre otros datos relevantes. Se recomienda no abusar de este tipo de fuentes.
\end{enumerate}




\section{Recomendaciones}
A la hora de abordar un trabajo como este, de cierta complejidad y extensión, es conveniente tener ciertas consideraciones desde un principio que ayuden a la organización y realización del mismo.
\begin{itemize}
    \item La memoria deberá ceñirse a las directrices dadas en la sección precedente. 
    
    \item Cualquier consulta externa (libro, artículo, página web, imagen,\ldots) debe estar debidamente referenciada tanto en el texto como en la bibliografía al final del trabajo. La bibliografía debe de aparecer en orden alfabético (del primer autor) en el formato indicado en la plantilla. 

    \item Se debe evitar copiar texto de forma literal, salvo citas literales, que se indicarán como tales y entrecomilladas. LaTeX proporciona el entorno \texttt{quote} para ello.

    \item Todas las imágenes y tablas incluidas en el documento deben figurar con su respectivos créditos (excepto que sean de elaboración propia). Por tanto, es recomendable guardar las referencias consultadas (direcciones web, libros) para la obtención de cualquier material gráfico o de datos.

    \item Si el trabajo contiene gran cantidad de vocabulario específico, conviene añadir un glosario de términos al final del mismo. Esto es mejor ir haciéndolo conforme se avanza en la redacción del trabajo.

    \item Es conveniente hacer un esquema inicial con la estructura general de la memoria: ¿de cuántas partes constará? ¿en qué orden? ¿qué incluirá cada una de ellas? En la plantilla proporcionada se recomienda una estructura general. Ello ayudará a organizar mejor el trabajo. No obstante, dicha estructura inicial puede ser modificada cuando el trabajo esté avanzado si el contenido lo requiere.
\end{itemize}


% Añadir tantos capítulos como sea necesario

\cleardoublepage\part{Segunda parte}
% !TeX root = ../tfg.tex
% !TeX encoding = utf8
%
%*******************************************************
% Introducción
%*******************************************************

% \manualmark
% \markboth{\textsc{Introducción}}{\textsc{Introducción}} 

\chapter{Introducción}

% !TeX root = ../tfg.tex
% !TeX encoding = utf8

\chapter{Preliminares}
Let x represent a discrete variable. As you learned in Chapter 3, the unit discrete
impulse, d(x), serves the same purposes in the context of discrete systems as the
impulse d(t) does when working with continuous variables. It is defined as


Teorema de nyquist.

bandlimited

uunit impulse 
226 sampling
contnuo
intervalo sampling theorem
\chapter{Transformada de Fourier Discreta}


\section{Introducción}

En este capítulo, analizaremos  en detalle la Transformada Discreta de Fourier (DFT), una herramienta crucial en el procesamiento digital de señales. Estudiaremos algunas de sus propiedades más importantes, que no deberán de sorprender al lector, ya que resultarán familiares y vendrán motivadas por su similitud con las propiedades descritas en el ámbito continuo. Haremos especial hincapié en el teorema de convolución cuya importancia quedará demostrada en el último capítulo del presente trabajo.

Se podría presentar directamente la definición de la DFT sin necesidad de justificar su origen. Sin embargo se opta por tratar de desarrollar la intuición de la definición de esta, ya que proporcionar una conexión conceptual con el dominio continuo contribuye a una comprensión más profunda y coherente de la Transformada Discreta de Fourier y de sus propiedades.

Es interesante mencionar que la mayoría de aproximaciones de la DFT tienen su origen en las Series de Fourier, y puede el lector preguntarse porque hemos establecido un marco teórico distinto del que partir. La razón es que conocer otras posibles interpretaciones, y tratar de relacionar la DFT con la Transformada de Fourier continua puede aportar una comprensión más profunda y completa de esta herramienta. 
Por ejemplo en visión por computador aparece, entre otros, el concepto de  ``derivación de la convolución'', que es el ``análogo'' al estudiado en la primera parte del trabajo, y por tanto proporcionan una intuición directa del mismo.

\vspace{0.2mm}

\noindent Puede por tanto el lector pensar que estamos tratando de transformar lo ``continuo'' en ``discreto'' en la medida de lo posible.

\section{Transformada de Fourier discreta}

Sin perder de vista nuestro objetivo, queremos calcular la Transformada de Fourier de una determinada imagen, de tal manera que deberíamos estar interesados en desarrollar una expresión 2-D que conserve las propiedades naturales que cabría esperar. Sin embargo de ahora en adelante prestaremos especial atención a la transformada 1D, cuyas propiedades y eficiencia computacional estudiaremos al detalle.


Esto se debe a que, como hemos mencionado previamente, la Transformada de Fourier discreta 2D, puede expresarse en función de la  DFT 1D. Tanto es así que de hecho, las librerías que calculan la DFT (usando el algoritmo FFT, descrito en el siguiente capítulo) de una determinada señal 2D, en realidad aplican sucesivas transformadas 1D y por tanto usan el algoritmo FFT 1-D. Por esta razón, a partir de este punto, nos concentraremos en el estudio de la DFT 1D, sin perder de vista ciertos aspectos relevantes de la DFT 2D.

\subsection{Transformada de Fourier discreta 1D}

Sea $f : \mathbb{R} \rightarrow  \mathbb{R}$  una función integrable en $\mathbb{R}$ que representará una señal 1D.

\noindent Haremos por tanto dos suposiciones acerca de $f$.
\begin{itemize}
    \item Suponemos que \( f \) es de soporte compacto en \([0, L]\), donde \( L \in \mathbb{N} \). Es decir $f(t)=0 \, \, \forall t \in \mathbb{R} \setminus [0, L]$.
    \item Suponemos que \( \widehat{f} \) es de soporte compacto en  \([0, 2B]\), donde \( B \in \mathbb{N} \). Esto es $\widehat{f}(y)=0 \, \, \forall y \in \mathbb{R} \setminus [0, 2B]$.
\end{itemize}

\noindent Estamos considerando  por tanto que $f$ es tanto limitada en el tiempo como limitada en banda, sabiendo que esto no tiene no es en general cierto. Sin embargo, finalmente vamos a llegar a una definición de una transformada discreta de Fourier que tendrá sentido por sí misma, independientemente de estas precarias suposiciones iniciales, que por otra parte tienen sentido al conocer las limitaciones computacionales que nos llevan a trabajar con operaciones finitas.  

\begin{itemize}
    \item \textbf{Paso 1} \\
    Tomamos $N$ muestras equiespaciadas de $f$  en $[0,L]$ con \[ N = \frac{L}{\frac{1}{2B}} = 2BL. \]

De manera que la tasa de muestreo será de $2B$ ejemplos por segundo. Obteniendo  los puntos $t_k$:
\begin{equation}
   t_0 = 0, \quad t_1 = \frac{1}{2B}, \quad t_2 = \frac{2}{2B}, \ldots, t_{N-1} = \frac{N-1}{2B}. 
\end{equation}

Consideramos la versión discreta de \( f \) como la lista \( f_d = [f(t_0), f(t_1), \ldots, f(t_{N-1})] \). 
\begin{observacion}
 Lejos de ser arbitraria, esta tasa de muestreo que hemos tomado viene inspirada en el Teorema de muestreo de Nyquist-Shannon que asegura que la construcción exacta de una señal continua en banda base a partir de sus muestras, es matemáticamente posible, si la señal está limitada en banda y la tasa de muestreo es superior al doble de su ancho de banda. En este caso la estamos considerando exactamente el doble de la frecuencia. Muestrear exactamente a la tasa de Nyquist a veces es suficiente para una recuperación perfecta de la función, pero hay casos en los que esto conduce a ciertas dificultades [ref].
\end{observacion}


\item \textbf{Paso 2} \\ 
Podemos expresar $f_d : \mathbb{R} \rightarrow \mathbb{R}$ de la siguiente forma
\begin{equation}
    f_d(t) = \sum_{k=0}^{N-1} \delta(t - t_k)  f(t_k) \, \,, t \in \mathbb{R}.
\end{equation} 
Observamos que $f_d$ corresponde con la función $f$ discretizada.  Para su construcción se ha usado la función impulso discreta $\delta$ [ref].

Se tiene evidentemente que $f_d$ es integrable en $\mathbb{R}$. Construimos por tanto $\widehat{f_d}$ usando la definición \ref{def:definicion} para $n=1$, y obtenemos  $\widehat{f_d}: \mathbb{R}  \rightarrow \mathbb{C}$
\begin{equation}\label{eqd}
    \widehat{f_d}(s) = \sum_{k=0}^{N-1} f(t_k)   \widehat{\delta}(t - t_k) = \sum_{n=0}^{N-1} \int_{\mathbb{R}}f(t_k)   \delta(t - t_k)e^{-2 \pi i t s} \, dt = \sum_{n=0}^{N-1} f(t_k)  e^{-2\pi i s t_k}.
\end{equation}

\item  \textbf{Paso 3}\\
Hemos calculado la Transformada de Fourier continua de la función discreta $f_d$, obteniendo una función $\widehat{f_d} \,$ definida en $\mathbb{R}$ que requiere ser discretizada. Recordamos que por hipótesis sabemos que $\widehat{f_d}(t) = 0 \, \, \forall t \notin [0,2B]$. Siguiendo con la filosofía usada anteriormente, usamos una tasa de muestreo de $\frac{1}{L}$, tomamos entonces $N$ muestras equiespaciadas de $\widehat{f_d}$   con 
\begin{equation}
    N = \frac{2B}{\frac{1}{L}} = 2BL.
\end{equation}
Obteniendo  los puntos $s_m$
\begin{equation}
    s_0 = 0, \quad s_1 = \frac{1}{L}, \quad s_2 = \frac{2}{L}, \ldots, s_{N-1} = \frac{N-1}{L}. 
\end{equation}
Consideramos por tanto la versión discreta de $\widehat{g}$ como la lista \( \, \widehat{g}_d =[\widehat{g} (s_0), \widehat{g} (s_1), \ldots, \widehat{g} (s_{N-1})] \). Aplicando~\eqref{eqd} se tiene finalmente que 
\begin{equation}\label{eq:tras}
    \widehat{f_d}(s_m) = \sum_{k=0}^{N-1} f(t_k) e^{-2\pi i s_m t_k}, \quad m = 0,1,\ldots,N-1
\end{equation} 
que nos da una expresión para pasar de la versión discreta de $f_d$ a la versión discreta de su transformada $\widehat{f_d}$.
\end{itemize}
\noindent Cabe preguntarse en qué sentido esta expresión es una aproximación discreta de la transformada de Fourier de $f$. Esta cuestión se puede abordar de la siguiente manera.

\noindent Calculamos la Transformada de Fourier de $f$;
\begin{equation}
  \widehat{f}(s)= \int_{0}^{L} e^{-2\pi ist} f(t) \, dt.  
\end{equation}
Evaluándola en los puntos de muestra $s_m$, se tiene
\begin{equation}
   \widehat{f}(s_m) = \int_{0}^{L} e^{-2\pi is_m t} f(t) \, dt \quad  \forall m \in \{0, \ldots,N-1\}
\end{equation}
Empleando ahora la aproximación de suma de Riemann para la integral con los puntos $t_n$, los cuales estaban a distancia  $\Delta t = \frac{1}{2B}$,  se tiene que dado $m \in \{0, \ldots,N-1\}$
\begin{equation} \label{eq:di}
\begin{aligned}
\widehat{f}(s_m) &= \int_{0}^{L} f(t)e^{-2\pi ism t} \, dt \\
&\approx  \sum_{n=0}^{N-1} f(t_n) e^{-2\pi is_m t_n} \Delta t = \frac{1}{2B} \sum_{n=0}^{N-1} f(t_n) e^{-2\pi is_m t_n} = \frac{1}{2B} \widehat{f_d}(s_m).
\end{aligned}
\end{equation}


\noindent Si atendemos a~\eqref{eq:di}, los valores $\widehat{f_d}(s_m)$ proporcionan una aproximación de los valores $\widehat{f}(s_m)$ salvo una cierta constante  $\frac{1}{2B}$.

\noindent Si usamos la definición de los puntos $t_k = \frac{k}{2B}$ y $s_m = \frac{m}{L}$ y reescribimos~\eqref{eq:tras} se tiene que 
\begin{equation}\label{eq:deft}
     \widehat{f_d}(s_m) = \sum_{k=0}^{N-1} f(t_k) e^{-2\pi i s_n t_k} = \sum_{k=0}^{N-1} f(t_k) e^{-2\pi i \frac{nm}{2BL}} = \sum_{k=0}^{N-1} f(t_k) e^{-2\pi i \frac{nm}{N}}.
\end{equation}
Lo verdaderamente notable de la ecuación \eqref{eq:deft} radica en que el factor $e^{-2\pi i \frac{nm}{N}}$ depende únicamente de las constantes $n$, $m$ y $N$, las cuales representan respectivamente el número de muestras de entrada, de salida y el número total de puntos. Este factor oculta los valores de las muestras en sí mismos. En cuanto al otro factor $f(t_k)$, iteramos sobre el índice k, y es en realidad a partir de este índice como recuperamos el valor correspondiente. Estas observaciones nos llevan a una definición más general que tiene un significado intrínseco. En lugar de concebir la DFT en términos de los valores muestreados de una señal continua y el muestreo de su transformada, la entenderemos como una operación que toma una lista con N elementos y produce otra lista con N elementos. 

\noindent Presentamos finalmente la definición de la Transformada de Fourier Discreta en 1D.

\begin{definicion}  \label{def:tfd} 
Sea $f = (f[0],f[1],\ldots,f[N-1])$ una N-tupla. Definimos la Transformada de Fourier Discreta de $f$ como la N-tupla $\widehat{f} = (F[0],F[1],\ldots,F[N-1])$ definida como:
\begin{equation}
    \widehat{f}[m] = \sum_{k=0}^{N-1} f[k] e^{-2\pi i \frac{mk}{N}}, \quad m = 0,1,\ldots,N-1
\end{equation}
\end{definicion}


\noindent La DFT define una biyección en el espacio $\mathbb{C}^N$. En efecto la transformación $\,\widehat{} \,: \mathbb{C}^N \rightarrow \mathbb{C}^N$ que a un vector $y = (y_0, \ldots, y_n) \in \mathbb{C}^N$, le hace corresponder otro vector $Y = (Y_0, \ldots, Y_n) \in \mathbb{C}^N$ dado por:
\begin{equation}
    Y_n = \sum_{k=0}^{N-1} y_k e^{-2\pi i \frac{mn}{N}}, \quad m = 0,1,\ldots,N-1
\end{equation}
tiene como inversa 
\begin{equation}
    y_n = \frac{1}{N}\sum_{k=0}^{N-1} Y_k e^{2\pi i \frac{mn}{N}}, \quad m = 0,1,\ldots,N-1
\end{equation}

\begin{observacion}
    En lo sucesivo adoptaremos el siguiente convenio:
    Al operar con la DFT de una secuencia de $N$ elementos  $\{y_0, y_1, y_2, \ldots, y_{N-1}\} $, consideramos que dicho vector es una muestra de una sucesión infinita periódica con período \( N \). Por lo que para cualquier entero \( k \) arbitrario, \( y_k \) se define como \( y_q \), donde \( 0 \leq q \leq N - 1 \) es el residuo de la división de \( k \) por \( N \). De la misma manera el vector $Y$ es periódico,  \( Y_{k+N} = Y_k \). Por tanto la TFD transforma señales periódicas discretas en el dominio del tiempo en señales periódicas discretas en el dominio de la frecuencia.
\end{observacion}


\section{Propiedades}

\subsection{Linealidad}
trivial.
\subsection{Simetrías}

\subsection{Periodicidad}
hablar de la periodicidad
y tb 
\subsection{Inversión}
en realidad hecho
\subsection{Expresión Matricial} no necesario



\noindent De donde se sigue inmediatamente la siguiente definición.

\begin{definicion}   
Sea $F = (F[0],F[1],\ldots,F[N-1])$ una N-tupla. Definimos la Inversa de la Transformada de Fourier Discreta de $F$ como la N-tupla $f = (f[0],f[1],\ldots,f[N-1])$ definida como:
\begin{equation}
    f[n] = \frac{1}{N}\sum_{m=0}^{N-1} F[m] e^{2\pi i \frac{mn}{N}}, \quad n = 0,1,\ldots,N-1
\end{equation}
\end{definicion}

\noindent Recibe el nombre de Inversa de la Transformada de Fourier Discreta ya que a partir de una comprobación directa obtenemos la identidad al componer ambas.




\subsection{Transformada de Fourier discreta 2D}
El razonamiento anterior puede repetirse para deducir la expresión de la Transformada Discreta 2D. Sin embargo si atendemos a la expresión de la Transformada de Fourier en $R^2$, descrita en la primera parte.

\noindent Sea $f: \mathbb{R}^2 \rightarrow \mathbb{C}$, tal que $f$ sea integrable en $\mathbb{R}^2$, se tiene que dado $(y_1,y_2) \in \mathbb{R}^2$: 
\begin{equation}
    \widehat{f}(y_1,y_2) = \int_{\mathbb{R}^2} f(x_1,x_2) e^{-2\pi i (x_1y_1+x_2y_2)} \, dx_1dx_2   =  \int_{\mathbb{R}}e^{-2\pi i (x_2y_2)}\int_{\mathbb{R}} f(x_1,x_2) e^{-2\pi i (x_1y_1)} \, dx_1dx_2 
\end{equation}

\noindent Si hacemos un esfuerzo y pensamos  esta expresión de manera discreta, tenemos que para un valor concreto de $x_2$, movemos $x_1$ a lo largo de todo el dominio considerado,  calculando  $ f(x_1,x_2)e^{-2\pi i (x_2y_2)}$ y acumulando los resultados obtenidos . Este proceso se repite con todos los valores de $x_2$ en el dominio correspondiente.

\noindent En efecto, esto nos lleva a la definición de la Transformada de Fourier discreta 2D:

\begin{definicion}   
Sea $F = (f[n,m])$ una matriz de tamaño $N \times M$. Definimos la Transformada de Fourier Discreta bidimensional de $f$ como la matriz $\widehat{F} = (F[k,l])$ definida por:

\begin{equation}
     \widehat{F}[k,l] = \sum_{n=0}^{N-1} \sum_{m=0}^{M-1} f[n,m] e^{-2\pi i \left(\frac{kn}{N} + \frac{lm}{M}\right)}, \quad k = 0,1,\ldots,N-1 \quad \text{y} \quad l = 0,1,\ldots,M-1
\end{equation}
\end{definicion}


\noindent Dado $k \in \{0,1,\ldots,N-1\}$ y $l \in \{0,1,\ldots,M-1\}$. Veamos que podemos escribir la DFT como
 \begin{equation} \label{eq:c}
     \widehat{F}[k,l] = \sum_{n=0}^{N-1} e^{-2\pi i \left(\frac{kn}{N}\right) }\sum_{m=0}^{M-1} f[n,m] e^{-2\pi i \left(\frac{lm}{M}\right) } = \sum_{n=0}^{N-1} e^{-2\pi i \left(\frac{kn}{N}\right) } G[n,l]. 
\end{equation}
donde dado $n \in \{0,1,\ldots,N-1\}$
\begin{equation} \label{eq:c2}
    G[n,l]=\sum_{m=0}^{M-1} f[n,m] e^{-2\pi i \left(\frac{lm}{M}\right) }
\end{equation}


 
\noindent Para un valor de \(n\) observamos que \(G(n,l)\) es la DFT 1-D de la n-ésima fila de \(f(x, y)\). Al variar \(n\) desde 0 hasta \(M - 1\) en la ecuación~\eqref{eq:c2}, calculamos un conjunto de DFT 1-D para todas las filas de \(f(x, y)\). Posteriormente se le realizan transformaciones 1-D a las columnas de este conjunto resultante \(G(n,l)\). Así, concluimos que la DFT 2-D de \(f(x, y)\) puede obtenerse mediante el cálculo de la transformada 1-D de cada fila de \(f(x, y)\) y luego calcular la Transformada 1-D a lo largo de cada columna del resultado. Esta es una simplificación importante porque solo tenemos que lidiar con una variable a la vez.
Esta propiedad es conocida como \textit{separabilidad}.





\section{Convolución}

Distinguimos convolución circular






\subsection{Convolución 1D}

\subsection{Convolución 2D}
\begin{definicion}
Sean $f[n], g[n]$ dos secuencias de  $N$. Definimos la convolución circular de $f[n]$ con $g[n]$ como la secuencia $h[n]$ definida por:
\begin{equation}
     h[n] = \sum_{k=0}^{N-1} f[k] \cdot g[(n-k)], \quad n = 0,1,\ldots,N-1.
\end{equation}
\end{definicion}




\subsection{Teorema de Convolución}
\noindent

\noindent Abordamos finalmente el teorema de convolución, el cual nos proporciona una herramienta fundamental para implementar convoluciones como productos puntuales. Esto sugiere una alternativa al cómputo directo de la convolución, abriendo nuevas posibilidades en diversos ámbitos, incluyendo en el de visión por computador.

\noindent Este resultado, al igual que los anteriores, es idéntico al presentado en la primera parte, a pesar de que hayamos cambiado el marco conceptual en el que nos movemos.


\begin{teorema}{(de Convolución 1D.)}
Sea \( x(n) \) una secuencia de \( N \) elementos y sea \( h(n) \) una secuencia de \( M \) elementos. Entonces se tiene que  
 \[
\widehat{(X*H)}(k) = \widehat{X}(k) \cdot \widehat{H}(k)
\]

\end{teorema}

\begin{proof}
\end{proof}

\begin{teorema}{(de Convolución 2D.)}
Sea \( x(n) \) una secuencia de \( N \) elementos y sea \( h(n) \) una secuencia de \( M \) elementos. Entonces se tiene que  
 \[
\widehat{(X*H)}(k) = \widehat{X}(k) \cdot \widehat{H}(k)
\]

\end{teorema}

\begin{proof}
\end{proof}



ejemplo quizas con una imagen :-)

\subsection{Correlación}






























\endinput
%--------------------------------------------------------------------
% FIN DEL CAPÍTULO. 
%--------------------------------------------------------------------


\chapter{Transformada Rápida de Fourier}






\section{Motivación}
La Transformada Rápida de Fourier (FFT) es un algoritmo eficiente para calcular la Transformada de Fourier Discreta  y su inversa.

\noindent Su impacto en la era moderna es incalculable, prueba de ello es su inclusión por la revista IEEE de Computación en Ciencia e Ingeniería como uno de  los 10 algoritmos más influyentes en el desarrollo y la práctica de la ciencia y la ingeniería en el siglo XX [ref]


\noindent El capítulo comienza  contextualizando al lector con algunos referencias históricas del algoritmo.\noindent  Posteriormente avanzaremos con el desarrollo teórico que sustenta la FFT donde usaremos fundamentalmente la teoría desarrollada previamente. Este análisis  culminará con el estudio de la mejora en eficiencia en  comparación con el método directo y una implementación recursiva para un caso particular.


\noindent Para finalizar el capítulo se abordarán algunas de las aplicaciones en diversos campos, particularmente en el ámbito de visión por computador.Con esto se pretende convencer al lector de que la FFT se encuentra prácticamente en cualquier aplicación que involucre el análisis y procesamiento de señales.




\section{Historia}












Nos remontamos a 1945, año en el que finaliza la la Segunda Guerra Mundial tras el lanzamiento de las primeras bombas atómicas usadas directamente sobre una población civil en Japón por parte de Estados Unidos. Este acontecimiento desencadenó en  una imperiosa necesidad por parte de las naciones de controlar la proliferación nuclear al comprobar el grado de destrucción que podían provocar.

\noindent Como resultado, Estados Unidos celebró conversaciones con los soviéticos y otras naciones con capacidad nuclear con el objetivo final de detener la carrera de armamento nuclear.

\noindent Para hacerlo se trató de llegar a determinados acuerdos que prohibía la experimentación con energía nuclear.Sin embargo debido a la esperable desconfianza entre países, era necesario que cada nación tuviera la tecnología para identificar las pruebas nucleares realizadas por otros. 

\noindent Los hidrófonos podían detectarlas bajo el mar y los átomos residuales podían identificarse en el cielo desde pruebas en tierra. Pero detectar las pruebas subterraneas era todo un desafío.


\noindent Una idea era analizar series temporales
sismológicas obtenidas de sismómetros situados en distintos lugares estratégicos para estudiar estas señales y deducir si se trataba de un exprimento nuclear.
Dentro de una señal compleja, hay muchas ondas sinusoidales de diferentes frecuencias, todas contribuyendo a formar el resultado final, y esto puede descifrarse si podemos aislar cada frecuencia. Para ello  recurrieron a la DFT que permitía descomponer la señal en distintas frecuencias, sin embargo  debido a la gran complejidad computacional que se necesitaba este método era imposible de poner en práctica incluso con los ordenadores actuales.

\noindent No fue hasta 1960,cuando James Cooley y John Tukey  publicaron el algoritmo de FFT que reducía considerablemente el orden de complejidad de la DFT, pudiendo reducir el tiempo de cálculo de más de tres años a 35 minutos.

\noindent Trágicamente, era demasiado tarde para firmar una prohibición de pruebas completa, y las pruebas nucleares fueron forzadas bajo tierra, donde las pruebas aumentaron a una notable tasa de alrededor de una vez por semana.

\noindent Desde entonces, la FFT ha encontrado nuevos usos en casi todas las aplicaciones de comunicaciones y señalización de datos Sin embargo, tenía el potencial de ser mucho mayor: casi fue el algoritmo que detuvo la carrera armamentista nuclear, si hubiera llegado solo unos años antes.


\noindent Lo más fascinante es que la primera aparición de la Transformada Rápida de Fourier (FFT) data de 1807 del matemático Gauss.
Sus intereses estaban en ciertos cálculos astronómicos (un área recurrente de aplicación de la FFT) relacionados con la interpolación de órbitas asteroidales a partir de un conjunto finito de observaciones equidistantes. Seguramente, la perspectiva de un cálculo manual laborioso y extenso era una buena motivación para el desarrollo de un algoritmo rápido. Sin embargo su trabajo no fue publicado y apareció únicamente en sus obras recopiladas como un manuscrito inédito  [ref].

\noindent Algunos trabajos previos en los Cooley y Tukey se basaron se mencionan a continuación.

\noindent Un método eficiente para el cálculo de las interacciones de un experimento factorial $2^m$ fue introducido por Yates[1]. La generalización a $3^m$ fue proporcionada por Box et al. [1]. Good [2] generalizó estos métodos y proporcionó algoritmos elegantes, para los cuales una clase de aplicaciones es el cálculo de series de Fourier.

\noindent Otro precursor importante es el trabajo de Danielson y Lanczos, realizado en el servicio de la cristalografía de rayos X, otro usuario frecuente de la tecnología FFT. [ref]

\noindent A lo largo de la historia, estos y numerosos otros nombres han estado vinculados al algoritmo FFT culminado con Cooley y Tukey.
"Esto puede servir como motivación para explorar no solo enfoques novedosos, sino también para revisitar de vez en cuando viejos documentos y descubrir la diversidad de trucos e ideas ingeniosas que se empleaban en una época en la que el cálculo era una tarea laboriosa. Quizás entre las ideas descartadas antes de la era de las computadoras electrónicas, podamos encontrar valiosas semillas para nuevos algoritmos." [ref]

\begin{definicion}
    Orden de Complejidad, definir y decir clases de complejodad y llamarle T y presentar los ordenes de complejidad que vienen a continaución importante.
    Y difernecia entre complejidad del problema y del algoritmo y tal
\end{definicion}


\section{Método Directo}
Si atendemos a la expresión de la DFT.El enfoque directo implica utilizar el método de fuerza bruta, mediante el cual debemos realizar cada sumatoria para cada elemento de la N-tupla cuya transformada deseamos calcular.

\noindent De este modo, dada una N-tupla $x = [x_0,..,x_{N-1}]$ calcular la DFT  de $x$ implica para cada  $j \in \{0,...,N-1\}$ realizar $N$ operaciones.
Por tanto para calcular la salida completa se necesitarían $N$ operaciones por cada índice $j \in \{0,...,N-1\}$ lo que nos lleva a un total de $N^2$ operaciones.
Y por tanto tiene un orden de complejidad cuadrático $T = O(N^2)$
\begin{ejemplo}
\noindent Si disponemos de una imagen de $2048 \times 2048$ (tamaño habitual), esto significaría realizar del orden de 17 mil millones de operaciones para calcular una única 2DFT excluyendo las exponenciales que podrían calcularse una única vez y almacenarse.

\noindent  Mientras que calcular la FFT 2-D de una imagen de $2048 \times 2048$ requeriría del orden de 92 millones de operaciones, lo cual representa una reducción significativa respecto a los mil millones de cálculos mencionados anteriormente. 
\begin{observacion}
  \noindent En ambos casos se ha tenido en cuenta la factorización de una 2DFT en 1DFT para realizar una comparación justa.[insertar imagen grande ¿?]
\end{observacion}
\end{ejemplo}

\begin{observacion}
    Una operación consiste en una multiplicación y una suma de dos números complejos.
\end{observacion}

\section{FFT}
Sea $W_N = e^{\frac{-2\pi i}{N}}$ la $N$-ésima raíz de la unidad. 
La 1-DFT de $f = (f[0],f[1],\ldots,f[N-1])$ puede reescribirse como $F = (F[0],F[1],\ldots,F[N-1])$ donde:
\begin{equation}\label{}
    F[j] = \sum_{k=0}^{N-1} f[k] W_{N}^{jk}, \quad j \in \{0,1,\ldots,N-1\}
\end{equation}


\begin{teorema}
Sea  $N$  compuesto con factorización $N = r_1 \cdot r_2$ con $r_1,r_2$ divisores naturales no triviales de $N$.
La DFT de una $N$-tupla $f$  puede ser calculada en $T=O(N(r_1+r_2))$ 
\end{teorema}


\vspace{0.2cm}


\begin{proof}
    \noindent Reescribimos los índices $j,k$ como sigue:
\begin{equation*}
    j = j_1r_1+j_0  \quad  j_0 \in \{0,1,\ldots r_1-1\}, j_1 \in \{0,1,\ldots r_2-1\}
\end{equation*}
\begin{equation*}
     k = k_1r_1+k_0   \quad  k_0 \in \{0,1,\ldots r_2-1\}, k_1 \in \{0,1,\ldots r_1-1\}
\end{equation*}

\noindent Esto nos permite escribir la DFT en función de $j_0,j_1,k_0,k_1$.

\begin{align}
    F[(j_1,j_0)] &= \sum_{k_0=0}^{r_2-1}  \sum_{k_1=0}^{r_1-1} f(k_1,k_0) W_N^{(j_1r_1+j_0)(k_1r_2+k_0)} \nonumber \\
    &= \sum_{k_0=0}^{r_2-1}  \sum_{k_1=0}^{r_1-1} f(k_1,k_0) W_N^{(j_1r_1+j_0)k_1r_2}W_N^{(j_1r_1+j_0)k_0}  \nonumber \\
    &= \sum_{k_0=0}^{r_2-1}  \sum_{k_1=0}^{r_1-1} f(k_1,k_0) W_N^{j_0k_1r_2}W_N^{(j_1r_1+j_0)(k_0)}  \label{eq:fft3}
\end{align}

\noindent La última igualdad se justifica con:

\begin{equation}
    W_N^{(j_1r_1+j_0)k_1r_2} =  W_N^{j_1k_1r_1r_2}W_N^{j_0k_1r_2}=
    W_N^{j_1k_1N}W_N^{j_0k_1r_2}=W_N^{j_0k_1r_2}
\end{equation}
\noindent donde se ha usado que $W_N$ es la raíz enésima de la unidad.
Notamos que $f(k_1,k_0) W_N^{j_0k_1r_2}$ depende únicamente de $j_0$ y no de $k_0$. Esto nos permite definir la N-tupla $f_1$ como 
\begin{equation}\label{eq:fft1}
    f_1(j_0,k_0) = \sum_{k_1=0}^{r_1-1} f(k_1,k_0) W_N^{j_0k_1r_2}
\end{equation}
y podemos escribir ~\eqref{eq:fft3} en función de $f_1$

\begin{equation}\label{eq:fft2}
    F[(j_1,j_0)] = \sum_{k_0=0}^{r_2-1}  f_1(j_0,k_0)W_N^{(j_1r_1+j_0)(k_0)} 
\end{equation}

\noindent Observamos que reescribiendo los índices hemos podido desacoplar las dos sumatorias. Esta es la clave del cálculo de la FFT.


\indent Como $f_1$ es una $N$-tupla y cada elemento se calcula usando $r_1$ operaciones,necesitamos por un lado $Nr_1$ operaciones para el cáculo de $f_1$. Una vez obtenido $f_1$, $F$ es calculado en $r_2$ operaciones por cada elemento lo que nos deja un total de $Nr_2$ operaciones.
En total el algoritmo regido por las ecuaciones  ~\eqref{eq:fft1}, ~\eqref{eq:fft2} es del orden de $O(N(r_1+r_2))$
\end{proof}


\begin{teorema}
Sea $N$ compuesto con factorización $N = \gamma_1 \cdot \gamma_2 \cdot \ldots \cdot r_m$ con $\gamma_1, \gamma_2, \ldots, \gamma_m$ divisores naturales no triviales de $N$.
La DFT de una $N$-tupla $f$ puede ser calculada en $T = O(N(\gamma_1 + \gamma_2 + \ldots + \gamma_m))$.
\end{teorema}

\begin{proof}

La demostración consiste en un razonamiento recursivo. Se tiene que si $N = \gamma_1 \cdot \gamma_2 \cdot \ldots \cdot \gamma_m$ con $\gamma_1, \gamma_2, \ldots, \gamma_m$ dando $m$ pasos en el algoritmo anteriormente descrito se deduce que  $T=O(N(\gamma_1 + \gamma_2 + \ldots + \gamma_m))$.
\textbf{NO ESTÁ BIEN}.
Aplicaremos por tanto m veces el algoritmo descrito en [ref teorema anterior].Consideraremos en cada paso las variables $r_1$ y $r_2$ para una mayor claridad estas se irán renombrando en cada paso del algoritmo.Los índices $j_O,j_1,k_0,k_1$ serán los relativos al cálculo seguido en el caso anterior para la  factorización $N= r_1r_2$.

\vspace{0.1cm}
\boxed{Paso 1}

Sea  $ N_1 =  \underbrace{\gamma_1}_{r_1} \cdot \underbrace{\gamma_2 \cdot \ldots \cdot \gamma_m}_{r_2}$.

\begin{equation}
    F[(j_1,j_0)] = \sum_{k_0=0}^{r_2-1}  f_1(j_0,k_0)W_N^{(j_1r_1+j_0)(k_0)} = \sum_{k=0}^{r_2-1}  f_1(k)W_N^{jk} 
\end{equation}

donde $f_2$ es calculado usando $\gamma_1N$ operaciones mediante la siguiente expresión:
\begin{equation}
    f_1(j_0,k_0) = \sum_{k_1=0}^{r_1-1} f(k_1,k_0) W_N^{j_0k_1r_2}
\end{equation}
\vspace{0.1cm}
 
\boxed{Paso 2}
\vspace{0.1cm}

Sea $N_2 = \underbrace{\gamma_2}_{r_1} \cdot \underbrace{ \gamma_3 \cdot \ldots \cdot \gamma_m}_{r_2}$.
Aplicamos el proceso de nuevo ahora para la ecuación 
\begin{equation}\label{eq:fft2}
    F[j] = \sum_{k=0}^{N_2-1}  f_1(k)W_N^{jk}  
\end{equation}
Lo que nos deja con 
\begin{equation}\label{eq:fft2}
      F[j_s0,j_1]  = \sum_{k_0=0}^{r_2-1}  f_2(j_0,k_0)W_N^{(j_1r_1+j_0)(k_0)} = \sum_{k=0}^{r_2-1}  f_2(k)W_N^{jk} 
\end{equation}
donde $f_2$ es calculado usando $\gamma_2N$ operaciones mediante la siguiente expresión:
\begin{equation}\label{eq:fft1}
    f_2(j_0,k_0) = \sum_{k_1=0}^{r_1-1} f(k_1,k_0) W_N^{j_0k_1r_2}
\end{equation}


\boxed{Paso m}
\vspace{0.1cm}

Sea $N_m$ = $\underbrace{\gamma_{m-1}}_{r_1} \cdot \underbrace{\gamma_m}_{r_2}$.
Aplicamos el proceso de nuevo ahora para la ecuación 
\begin{equation}\label{eq:fft2}
   F[j] = \sum_{k=0}^{r_2-1}  f_m(k)W_N^{jk}  
\end{equation}
Lo que nos deja con 
\begin{equation}\label{eq:fft2}
    F[(j_1,j_0)] = \sum_{k_0=0}^{r_2-1}  f_2(j_0,k_0)W_N^{(j_1r_1+j_0)(k_0)} = \sum_{k=0}^{r_2-1}  f_1(k)W_N^{jk} 
\end{equation}

donde $f_m$ es calculado usando $\gamma_{m-1}N$ operaciones mediante la siguiente expresión: 
\begin{equation}\label{eq:fft1}
    f_m(j_0,k_0) = \sum_{k_1=0}^{r_1-1} f(k_1,k_0) W_N^{j_0k_1r_2}
\end{equation}
Finalmente concluimos con $r_mN$ operaciones el cálculo de $F$.

Lo que nos deja un total de $O(N(r_1+ \ldots r_m))$

Lo dejo apartado y ya está.





\end{proof}


\begin{observacion}
    Si queremos minimizar el orden de complejidad $T$, debemos usar tantos factores de $N$ como sea posible. Es decir si $N = r_1 \cdot r_2 \cdot \ldots \cdot r_m$ y dado $j \in \{1, \ldots m\}$ se tiene que $r_j = p_1p_2$ con $p_1,p_2>1$   entonces $p_1+p_2 < r_j$ salvo que  $p_1=p_2=2$.
    Y por consiguiente $O(N(r_1 + r_2 + \ldots + r_j + \ldots + r_m))$ tiene mayor orden de complejidad que $T=O(N(r_1 + r_2 + \ldots + p_1+p_2+ \ldots r_m))$.
    La igualdad se da  para $p_1=p_2=2$, de donde deducimos que el factor $2$ puede combinarse de manera indistinta sin pérdida alguna.
\end{observacion}


\textbf{hablar sobre la memoria tb}


\noindent Luego, al descomponer un número en factores, observamos una disminución del orden de complejidad. Por lo tanto, podemos inferir que los números con un mayor número de divisores proporcionarán una ganancia en eficiencia mayor.Esto motiva la siguiente definición:

\begin{definicion}
    Un número antiprimo o altamente compuesto es un entero positivo con más divisores que cualquier entero positivo más pequeño
\end{definicion}


\noindent Ramanujan (1915) enumeró 102 números altamente compuestos hasta 6746328388800, pero omitió algunos de ellos. Robin (1983) proporciona los primeros 5000 números altamente compuestos, y una encuesta exhaustiva se da por Nicholas (1988). Flammenkamp da una lista de los primeros 779674 números altamente compuestos.

\begin{ejemplo}
Son números antiprimos el 1, 2, 4, 6, 12, 24, 36, 48, 60.
\end{ejemplo}


\begin{proposicion}
  \noindent Si $N = 2^{a_2}3^{a_3} \ldots p^{a_p}$ es la factorización prima de un número altamente compuesto, \textbf{entonces}:
\end{proposicion}

\begin{itemize}
\item Los primos 2, 3, ..., p forman una secuencia de primos consecutivos,

\item Los exponentes son no decrecientes, por lo que \( a_2 \geq a_3 \geq \ldots \geq a_p \), y

\end{itemize}

\textbf{No lo se bien falta buscar mas datos sobre estos numeros}

\begin{proposicion}
    El algoritmo FFT proporciona ganancias significativas cuando el tamaño de entrada $N$ es un número altamente compuesto. 
\end{proposicion}

\begin{proof}
   El orden de complejidad de la FFT con una entrada de tamaño $N$ con  $N = p_1^{a_1}p_2^{a_2} \ldots p_m^{a_m}$ es $T= O(N(a_1\cdot p_1+a_2\cdot p_2+ \ldots a_m\cdot p_m$))

   \noindent Luego, 
    \begin{equation}
        \frac{T}{N} = a_1\cdot p_1+a_2\cdot p_2+ \ldots a_m\cdot p_m
    \end{equation}
\end{proof}

Aplicando log en base 2 \textbf{(por qué)}

\begin{equation}
    log_2N = log_2(a_1)\cdot p_1+log_2(a_2)\cdot p_2+ \ldots log_2(a_m)\cdot p_m
\end{equation}

De donde obtenemos
\begin{equation}
    \frac{T}{Nlog(N)} =  \frac{a_1\cdot p_1+a_2\cdot p_2+ \ldots a_m\cdot p_m}{log_2(a_1)\cdot p_1+log_2(a_2)\cdot p_2+ \ldots log_2(a_m)\cdot p_m}
\end{equation}

\noindent Cuando el número $N$ es altamente compositivo el término de la derecha se acota por una constante  $C$, y por tanto $\frac{T}{Nlog(N)} = C $, lo que implica que el orden de complejidad es de $T= O(Nlog(N)$. \textbf{esto deberíamos comprobarlo} \textbf{y quizás hablar sobre el tema de las ganancias}








\textbf{distinguie para los numeros para los que hay mas ganancia y hacer alguna grafica} comprobar porque no tengo claro lo de los números \textbf{}
Y grafica


\begin{proposicion}
 Si tomamos $r_1 = r_2 = \ldots = r_m = r$ entonces  $m = log_rN$ y  el número total de operaciones para el cálculo de la FFT con tamaño de entrada $N =r^m$  es $T(r) = rNlog_rN $
\end{proposicion}
\begin{proof}
Se deduce del teorema 5.2 [ref]

\end{proof}

\noindent El caso $r=2$ y $r=4$ es de especial interés a nivel computacional, ya que el algoritmo ofrece ciertas ventajas añadidas tanto en el direccionamiento como en las operaciones realizadas como consecuencia de la naturaleza binaria inherente a los ordenadores.

Detallaremos el caso $r = 2$. La lógica seguida durante el proceso es similar a la ya aplicada anteriormente.

Sea $N = 2^m$, expresamos los índices $j,k$ en [ref] como 
\begin{equation}
    j = j_{m-1}\cdot 2^{m-1}+ \ldots j_{1}\cdot 2 + j_0, \quad j_i \in \{0,1\} \quad \forall i \in \{0, \ldots m-1\}
\end{equation}
\begin{equation}
    k = k_{m-1}\cdot 2^{m-1}+ \ldots k_{1}\cdot 2 + k_0, \quad k_i \in \{0,1\} \quad \forall i \in \{0, \ldots m-1\}
\end{equation}
La clave está en que esta expresión de $j,k$ coincide con  la descomposición binaria de cada índice. De tal modo que $j_i, k_i$ refieren al bit que ocupa la posición $i$ de la representación binaria de $j$ y $k$ respectivamente, y por tanto solo pueden tomar el valor $0,1$.
Escribiremos entonces [ref def] en función de sus índices.Luego, \textbf{no lo tengo claro}

\begin{align}
    F[(j_1,j_0)] &= \sum_{k_0=0}^{1}  \sum_{k_1=0}^{1} \ldots \sum_{k_{m-1}=0}^{1} f(k_{m-1}, \ldots, k_0) W_N^{jk_{m-1}{2^{m-1}+ \ldots + jk_0}}
\end{align}

\noindent Siguiendo un razonamiento similar a [ref]. Como consecuencia de $W_N^N =1$.

\begin{equation}
    W_N^{jk_{m-1}{2^{m-1}}} =  W_N^{(j_{m-1}\cdot 2^{m-1}+ \ldots j_{1}\cdot 2 + j_0)k_{m-1}{2^{m-1}}}=W_N^{j_0k_{m-1}{2^{m-1}}}
\end{equation}

\noindent Esto nos permite  desacoplar la suma interior $f_1$ y calcularla directamtente  como 
\begin{equation}
    f_1 (j_0, k_{m-2}, \ldots k_0) = \sum_{k_{m-1}=0}^{1} f(k_{m-1}, \ldots, k_0) W_N^{j_0k_{m-1}{2^{m-1}}}
\end{equation}
ya que sólo depende de $j_0$.

\noindent Se procede de igual forma con el resto de índices desacoplando las distintas sumatorias.

\noindent De manera que en general para $l \in \{1, \ldots m\}$ tendríamos

\begin{equation}
    f_l (j_0, \ldots, j_{l-1},k_{m-l-1}, \ldots k_0) = \sum_{k_{m-l}=0}^{1} f_{l-1}(j_0, \ldots, j_{l-2},k_{m-l}, \ldots k_0)W_N^{(j_{l-1}2^{l-1}+\ldots+j_0)k_{m-l}2^{m-l}}
\end{equation}

\noindent donde hemos usado

\begin{equation}
    W_N^{jk_{m-l}{2^{m-l}}} =  W_N^{(j_{m-1}\cdot 2^{m-1}+ \ldots j_{1}\cdot 2 + j_0)k_{m-l}{2^{m-l}}}=W_N^{j_0k_{m-l}{2^{m-l}}}
\end{equation}
Podemos desarrollar la sumatoria y usar que $e^{i\pi}=-1$ $e^{i\frac{\pi}{2}=i}$
\begin{align*}
&f_l(j_0, \ldots, j_{l-1},k_{m-l-1}, \ldots, k_0) \\
&\quad = f_{l-1}(j_0, \ldots, j_{l-2},0,k_{m-l-1}, \ldots, k_0) \\
&\quad\quad + (-1)^{j_{l-1}}i^{j_{l-2}}f_{l-1}(j_0, \ldots, j_{l-2},1,k_{m-l-1}, \ldots, k_0)W_N^{(j_{l-3}2^{l-3}+\ldots+j_0)2^{m-l}} \quad j_{l-1} \in \{0,1\}  
\end{align*}

\noindent De acuerdo con la convención de indexación, esto se almacena en una ubicación cuyo índice es
\[ j_02^{m-1} + \cdot \cdot \cdot + j_{l-1}2^{m-l} + k_{m-l-1}2^{m-l-1} + \cdot \cdot \cdot + k_0 \]
Se puede observar en [ref] que solo están involucradas en el cálculo las dos ubicaciones de almacenamiento con índices que tienen $0$ y $1$ en la posición del bit \(2^{m-1}\). 

\noindent Además la computación paralela está permitida, ya que la operación descrita por [ref] se puede llevar a cabo con todos los valores de \(j_0, \cdot \cdot \cdot, j_{i-2}\) y \(k_0, \cdot \cdot \cdot, k_{m-i-1}\) simultáneamente. En algunas aplicaciones, es conveniente usar (20) para expresar \(A_i\) en términos de \(A_{i-2}\), lo que equivale a un  con \(r = 4\).

\noindent El último array $f_m$ calculado proporciona el resultado deseado de la siguiente manera.
\[ F(j_{m-i}, \ldots, j_0) = f_{m} (j_{0}, \ldots, j_{m-1}) \]




\section{Algoritmos}
Presentamos ahora algunos de los algoritmo más usados para la implementación de la transformada rápida de Fourier.

\textbf{RADIX-2:}

Implementación recursiva del  algoritmo FFT para el caso $N=2^m$.

\begin{algorithm}
\caption{FFT (Radix-2)}
\label{alg:fft}
\begin{algorithmic}[1]
\Procedure{FFT}{$A$}
    \State \textbf{Entrada:} Un array de valores complejos con tamaño $2^m$ para $m \geq 0$.
    \State \textbf{Salida:} Un array de valores complejos que es la DFT de la entrada.
    \State $N := \text{longitud}(A)$
    \If{$N = 1$}
        \State \textbf{devolver} $A$
    \Else
        \State $W_N := e^{\frac{2\pi i}{N}}$
        \State $W := 1$
        \State $A_{\text{par}} := (A_0, A_2, \ldots, A_{N-2})$
        \State $A_{\text{impar}} := (A_1, A_3, \ldots, A_{N-1})$
        \State $Y_{\text{par}} := \text{FFT}(A_{\text{par}})$
        \State $Y_{\text{impar}} := \text{FFT}(A_{\text{impar}})$
        \For{$j := 0$ \textbf{hasta} $\frac{N}{2} - 1$}
            \State $Y[j] = Y_{\text{par}}[j] + W \cdot Y_{\text{impar}}[j]$
            \State $Y[j + \frac{N}{2}] = Y_{\text{par}}[j] - W \cdot Y_{\text{impar}}[j]$
            \State $W := W \cdot W_N$
        \EndFor
        \State \textbf{devolver} $Y$
    \EndIf
\EndProcedure
\end{algorithmic}
\end{algorithm}



\textbf{RADIX-4:}

Implementación recursiva del  algoritmo FFT para el caso $N=4^m$.









\section{Rader}
El algoritmo de Cooley-Tukey de la FFT no contempla para $N$ primo ninguna mejora.
La mayoría de algoritmos de la FFT aplican en este caso un zero-padding  a la potencia de 2 más cercana para aproximar la transformada de fourier continua.
Sin embargo si lo que se quiere es calcular la DFT de una secuencia de elementos podemos recurrir al algoritmo de Rader que utiliza la teoría de números  para proporcionar una alternativa en este caso.


Recordamos que  $\mathbb{Z}_p = \mathbb{Z}/p$, esto es, el anillo cociente $\mathbb{Z}$ módulo el ideal  ${p\mathbb{Z}}$, que es un cuerpo gracias a que $p\mathbb{Z}$ es un ideal maximal.
Denotamos por $\mathbb{Z}_p^\times$ las unidades de 
$\mathbb{Z}_p.$

\begin{teorema}
        $\mathbb{Z}_p^\times$ es un grupo cíclico.
\end{teorema}

\begin{teorema}
    (Teorema de Fermat)
\end{teorema}

El algoritmo de Rader consiste por tanto en :
 



bla bla


primos de sophie germain y cadena ytal.
----------------------------------



 









\chapter{Aplicaciones}


\section{Motivación}
\textbf{quizas esto deberia ser otro capitulo o deberia estar al final de lo de DFT}
"Si aceleras cualquier algoritmo no trivial en un factor de un millón o más, el mundo encontrará rápidamente aplicaciones útiles para él " [ref]

\noindent Mostramos a continuación algunas de las posibles aplicaciones de la FFT en el campo de visión por computador.

\section{Cálculo de Polinomios}

\section{Tratamiento de Imágenes}

\begin{itemize}
    \item Reducción de Ruido
    \item Compresión de Imágenes
    \item Pegado de Imágenes

\end{itemize}

\section{Detección de Características}
\begin{itemize}
    \item Detección de Bordes
    \item Detección de puntos de interés
\end{itemize}

\section{Aplicaciones en Redes Neuronales}
El culmen del tfga
teorema de convolucion y ta,

% -------------------------------------------------------------------
% APPENDIX: Opcional
% -------------------------------------------------------------------

\appendix % Reinicia la numeración de los capítulos y usa letras para numerarlos
\pdfbookmark[-1]{Apéndices}{appendix} % Alternativamente podemos agrupar los apéndices con un nuevo \part{Apéndices}

% !TeX root = ../tfg.tex
% !TeX encoding = utf8

\chapter{Apéndices}\label{ap:apendice1}


\begin{itemize}
  \item Desigualdad de Young
  \item Desigualdad de Holder
  \item Fubbini y Tonelli
  \item 
\end{itemize}

\endinput
%------------------------------------------------------------------------------------
% FIN DEL APÉNDICE. 
%------------------------------------------------------------------------------------

% Añadir tantos apéndices como sea necesario 

% -------------------------------------------------------------------
% GLOSARIO: Opcional
% -------------------------------------------------------------------

% !TeX root = ../tfg.tex
% !TeX encoding = utf8

\chapter*{Glosario}
\addcontentsline{toc}{chapter}{Glosario} % Añade el glosario a la tabla de contenidos

La inclusión de un glosario es opcional.

Archivo: \texttt{glosario.tex}

\begin{description} 
  \item[$\mathbb{R}$] Conjunto de números reales.

  \item[$\mathbb{C}$] Conjunto de números complejos.

  \item[$\mathbb{Z}$] Conjunto de números enteros.
\end{description}
\endinput
 

% -------------------------------------------------------------------
% BACKMATTER
% -------------------------------------------------------------------

\backmatter % Desactiva la numeración de los capítulos
\pdfbookmark[-1]{Referencias}{BM-Referencias}

% BIBLIOGRAFÍA
%-------------------------------------------------------------------

\bibliographystyle{alpha-es} 
\begin{small} 
  \bibliography{library.bib}
\end{small}


\end{document}
